% El método de las soluciones manufacturadas, paso a paso.
% Todas las cifras vienen de datos.tex (generar_datos.py) y de
% datos_sabotaje.tex (experimento_sabotaje.py). Compilar:
%   pdflatex mms_paso_a_paso.tex   (dos veces, por las referencias cruzadas)
\documentclass[11pt,a4paper]{article}
\usepackage[utf8]{inputenc}
\usepackage[T1]{fontenc}
\usepackage{lmodern}
\usepackage[spanish,es-tabla,es-noshorthands]{babel}
\usepackage[a4paper,margin=2.4cm]{geometry}
\usepackage{microtype}
\usepackage{amsmath,amssymb,mathtools,bm}
\usepackage{icomma}
\usepackage{booktabs,array,multirow}
\usepackage[table]{xcolor}
\usepackage{tikz}
\usetikzlibrary{arrows.meta,positioning,calc,babel}
\usepackage{pgfplots}
\pgfplotsset{compat=1.18}
\usepgfplotslibrary{groupplots}
\usepackage{caption}
\usepackage{float}
\usepackage{enumitem}
\usepackage[skins,breakable]{tcolorbox}
\usepackage{listings}
\usepackage{pifont}
\usepackage[hidelinks]{hyperref}
\hypersetup{pdftitle={El método de las soluciones manufacturadas, paso a paso},
            pdfauthor={EduFEM}, pdfsubject={Verificación del motor de EduFEM}}

% --- Paleta (referencia de la guía de visualización) ------------------------------
\definecolor{tinta}{HTML}{0B0B0B}
\definecolor{tintasec}{HTML}{52514E}
\definecolor{apagado}{HTML}{898781}
\definecolor{grilla}{HTML}{E1E0D9}
\definecolor{ejeC}{HTML}{C3C2B7}
\definecolor{fondo}{HTML}{F4F3EF}
\definecolor{serieA}{HTML}{2A78D6}  % Q4
\definecolor{serieB}{HTML}{EB6834}  % Q9
\definecolor{bien}{HTML}{0CA30C}
\definecolor{mal}{HTML}{D03B3B}

\setlength{\parindent}{0pt}
\setlength{\parskip}{0.55em}
\captionsetup{font=small, labelfont=bf, margin=0.5cm}
\setlist{itemsep=0.2em, topsep=0.3em}

% --- Datos generados ----------------------------------------------------------------
\newcommand{\dato}[1]{\ifcsname dato:#1\endcsname\csname dato:#1\endcsname\else\textbf{\textcolor{mal}{??#1}}\fi}
\newcommand{\tabla}[1]{\csname tabla:#1\endcsname}
\newcommand{\cmark}{\textcolor{bien}{\ding{51}}}
\newcommand{\xmark}{\textcolor{mal}{\ding{55}}}
% Generado por generar_datos.py: no editar a mano.
\expandafter\def\csname dato:E\endcsname{1{,}0}
\expandafter\def\csname dato:nu\endcsname{0{,}3}
\expandafter\def\csname dato:t\endcsname{1{,}0}
\expandafter\def\csname dato:c\endcsname{1{,}098901}
\expandafter\def\csname dato:mu\endcsname{0{,}384615}
\expandafter\def\csname dato:dosmu\endcsname{0{,}769231}
\expandafter\def\csname dato:lam\endcsname{0{,}329670}
\expandafter\def\csname dato:K2\endcsname{7{,}592003}
\expandafter\def\csname dato:pi\endcsname{3{,}141593}
\expandafter\def\csname dato:pi2\endcsname{9{,}869604}
\expandafter\def\csname dato:D11\endcsname{1{,}098901}
\expandafter\def\csname dato:D12\endcsname{0{,}329670}
\expandafter\def\csname dato:D33\endcsname{0{,}384615}
\expandafter\def\csname dato:DdifD\endcsname{0{,}769231}
\expandafter\def\csname dato:bmuestra\endcsname{3{,}796002}
\expandafter\def\csname dato:J11\endcsname{0{,}25}
\expandafter\def\csname dato:detJ\endcsname{0{,}0625}
\expandafter\def\csname dato:Jinv\endcsname{4}
\expandafter\def\csname dato:gdos\endcsname{0{,}577350}
\expandafter\def\csname dato:k1\endcsname{0{,}4500}
\expandafter\def\csname dato:ck1\endcsname{0{,}494505}
\expandafter\def\csname dato:k2\endcsname{0{,}1625}
\expandafter\def\csname dato:ck2\endcsname{0{,}178571}
\expandafter\def\csname dato:k3\endcsname{-0{,}2750}
\expandafter\def\csname dato:ck3\endcsname{-0{,}302198}
\expandafter\def\csname dato:k4\endcsname{-0{,}0125}
\expandafter\def\csname dato:ck4\endcsname{-0{,}013736}
\expandafter\def\csname dato:k5\endcsname{-0{,}2250}
\expandafter\def\csname dato:ck5\endcsname{-0{,}247253}
\expandafter\def\csname dato:k6\endcsname{-0{,}1625}
\expandafter\def\csname dato:ck6\endcsname{-0{,}178571}
\expandafter\def\csname dato:k7\endcsname{0{,}0500}
\expandafter\def\csname dato:ck7\endcsname{0{,}054945}
\expandafter\def\csname dato:k8\endcsname{0{,}0125}
\expandafter\def\csname dato:ck8\endcsname{0{,}013736}
\expandafter\def\csname dato:Bdiag\endcsname{-1{,}577350}
\expandafter\def\csname dato:Kcc\endcsname{1{,}978022}
\expandafter\def\csname dato:Kccxy\endcsname{0{,}000000}
\expandafter\def\csname dato:Kfrx\endcsname{-0{,}714286}
\expandafter\def\csname dato:Kfry\endcsname{0{,}000000}
\expandafter\def\csname dato:nesq\endcsname{0{,}178571}
\expandafter\def\csname dato:Kvecx\endcsname{0{,}109890}
\expandafter\def\csname dato:Kvecy\endcsname{-0{,}604396}
\expandafter\def\csname dato:gA\endcsname{0{,}105662}
\expandafter\def\csname dato:gB\endcsname{0{,}394338}
\expandafter\def\csname dato:sA\endcsname{0{,}325886}
\expandafter\def\csname dato:sB\endcsname{0{,}945409}
\expandafter\def\csname dato:Ncerca\endcsname{0{,}622008}
\expandafter\def\csname dato:Nmixto\endcsname{0{,}166667}
\expandafter\def\csname dato:Nlejos\endcsname{0{,}044658}
\expandafter\def\csname dato:sumaNss\endcsname{0{,}663391}
\expandafter\def\csname dato:fe1x\endcsname{0{,}314779}
\expandafter\def\csname dato:fe1y\endcsname{0{,}099015}
\expandafter\def\csname dato:F9\endcsname{1{,}259118}
\expandafter\def\csname dato:F10\endcsname{0{,}000000}
\expandafter\def\csname dato:Fredx\endcsname{1{,}973403}
\expandafter\def\csname dato:Fredy\endcsname{0{,}000000}
\expandafter\def\csname dato:u5\endcsname{0{,}997665}
\expandafter\def\csname dato:v5\endcsname{0{,}000000}
\expandafter\def\csname dato:erru5\endcsname{-0{,}002335}
\expandafter\def\csname dato:erru5pct\endcsname{0{,}23}
\expandafter\def\csname dato:sumaFx\endcsname{3{,}067531}
\expandafter\def\csname dato:sumaFy\endcsname{0{,}000000}
\expandafter\def\csname dato:sumaRx\endcsname{-3{,}067531}
\expandafter\def\csname dato:sumaRy\endcsname{0{,}000000}
\expandafter\def\csname dato:cx\endcsname{0{,}25}
\expandafter\def\csname dato:cy\endcsname{0{,}25}
\expandafter\def\csname dato:gux\endcsname{0{,}997665}
\expandafter\def\csname dato:guy\endcsname{0{,}997665}
\expandafter\def\csname dato:gvx\endcsname{-1{,}000000}
\expandafter\def\csname dato:gvy\endcsname{-1{,}000000}
\expandafter\def\csname dato:eux\endcsname{1{,}570796}
\expandafter\def\csname dato:euy\endcsname{1{,}570796}
\expandafter\def\csname dato:evx\endcsname{-1{,}570796}
\expandafter\def\csname dato:evy\endcsname{-1{,}570796}
\expandafter\def\csname dato:dux\endcsname{-0{,}573131}
\expandafter\def\csname dato:dvx\endcsname{0{,}570796}
\expandafter\def\csname dato:e2c\endcsname{1{,}308576}
\expandafter\def\csname dato:wc\endcsname{0{,}790123}
\expandafter\def\csname dato:dAc\endcsname{0{,}049383}
\expandafter\def\csname dato:cH1c\endcsname{0{,}064621}
\expandafter\def\csname dato:uhc\endcsname{0{,}249416}
\expandafter\def\csname dato:uMc\endcsname{0{,}500000}
\expandafter\def\csname dato:vhc\endcsname{0{,}250000}
\expandafter\def\csname dato:vMc\endcsname{0{,}500000}
\expandafter\def\csname dato:d2c\endcsname{0{,}125292}
\expandafter\def\csname dato:cL2c\endcsname{0{,}006187}
\expandafter\def\csname dato:sL2e1\endcsname{0{,}016376}
\expandafter\def\csname dato:sH1e1\endcsname{0{,}559242}
\expandafter\def\csname dato:totL2\endcsname{0{,}065503}
\expandafter\def\csname dato:totH1\endcsname{2{,}236968}
\expandafter\def\csname dato:totEx\endcsname{0{,}500000}
\expandafter\def\csname dato:L2\endcsname{0{,}255936}
\expandafter\def\csname dato:H1\endcsname{1{,}495650}
\expandafter\def\csname dato:normaEx\endcsname{0{,}707107}
\expandafter\def\csname dato:L2rel\endcsname{36{,}19}
\expandafter\def\csname dato:eL2dieciseis\endcsname{4{,}698601\times10^{-3}}
\expandafter\def\csname dato:eL2treintaydos\endcsname{1{,}176827\times10^{-3}}
\expandafter\def\csname dato:cocienteL2\endcsname{3{,}992600}
\expandafter\def\csname dato:tasaL2\endcsname{1{,}9973}
\expandafter\def\csname dato:tasaQcuatroLdos\endcsname{1{,}997}
\expandafter\def\csname dato:tasaQcuatroHuno\endcsname{1{,}001}
\expandafter\def\csname dato:tasaQnueveLdos\endcsname{3{,}000}
\expandafter\def\csname dato:tasaQnueveHuno\endcsname{2{,}001}
\expandafter\def\csname dato:tasaPrimeraQcuatro\endcsname{1{,}822}
\expandafter\def\csname dato:dosmupi\endcsname{2{,}416610}
\expandafter\def\csname dato:lamDP\endcsname{0{,}576923}
\expandafter\def\csname dato:sL2e1raiz\endcsname{0{,}127968}
\expandafter\def\csname dato:contraste\endcsname{coincide}
\expandafter\def\csname tabla:ke\endcsname{%
$u_1$ & $0{,}494505$ & $0{,}178571$ & $-0{,}302198$ & $-0{,}013736$ & $-0{,}247253$ & $-0{,}178571$ & $0{,}054945$ & $0{,}013736$ \\
$v_1$ & $0{,}178571$ & $0{,}494505$ & $0{,}013736$ & $0{,}054945$ & $-0{,}178571$ & $-0{,}247253$ & $-0{,}013736$ & $-0{,}302198$ \\
$u_2$ & $-0{,}302198$ & $0{,}013736$ & $0{,}494505$ & $-0{,}178571$ & $0{,}054945$ & $-0{,}013736$ & $-0{,}247253$ & $0{,}178571$ \\
$v_2$ & $-0{,}013736$ & $0{,}054945$ & $-0{,}178571$ & $0{,}494505$ & $0{,}013736$ & $-0{,}302198$ & $0{,}178571$ & $-0{,}247253$ \\
$u_3$ & $-0{,}247253$ & $-0{,}178571$ & $0{,}054945$ & $0{,}013736$ & $0{,}494505$ & $0{,}178571$ & $-0{,}302198$ & $-0{,}013736$ \\
$v_3$ & $-0{,}178571$ & $-0{,}247253$ & $-0{,}013736$ & $-0{,}302198$ & $0{,}178571$ & $0{,}494505$ & $0{,}013736$ & $0{,}054945$ \\
$u_4$ & $0{,}054945$ & $-0{,}013736$ & $-0{,}247253$ & $0{,}178571$ & $-0{,}302198$ & $0{,}013736$ & $0{,}494505$ & $-0{,}178571$ \\
$v_4$ & $0{,}013736$ & $-0{,}302198$ & $0{,}178571$ & $-0{,}247253$ & $-0{,}013736$ & $0{,}054945$ & $-0{,}178571$ & $0{,}494505$ \\
}
\expandafter\def\csname tabla:B\endcsname{%
$-1{,}577350$ & $0{,}000000$ & $1{,}577350$ & $0{,}000000$ & $0{,}422650$ & $0{,}000000$ & $-0{,}422650$ & $0{,}000000$ \\
$0{,}000000$ & $-1{,}577350$ & $0{,}000000$ & $-0{,}422650$ & $0{,}000000$ & $0{,}422650$ & $0{,}000000$ & $1{,}577350$ \\
$-1{,}577350$ & $-1{,}577350$ & $-0{,}422650$ & $1{,}577350$ & $0{,}422650$ & $0{,}422650$ & $1{,}577350$ & $-0{,}422650$ \\
}
\expandafter\def\csname tabla:nodos\endcsname{%
1 & $0{,}0$ & $0{,}0$ & borde & $0$ & $1$ \\
2 & $0{,}5$ & $0{,}0$ & borde & $0$ & $0$ \\
3 & $1{,}0$ & $0{,}0$ & borde & $0$ & $-1$ \\
4 & $0{,}0$ & $0{,}5$ & borde & $0$ & $0$ \\
5 & $0{,}5$ & $0{,}5$ & libre & \multicolumn{2}{c}{incógnitas $u_5,\ v_5$} \\
6 & $1{,}0$ & $0{,}5$ & borde & $0$ & $0$ \\
7 & $0{,}0$ & $1{,}0$ & borde & $0$ & $-1$ \\
8 & $0{,}5$ & $1{,}0$ & borde & $0$ & $0$ \\
9 & $1{,}0$ & $1{,}0$ & borde & $0$ & $1$ \\
}
\expandafter\def\csname tabla:Kdiag\endcsname{%
$e_1$ & 3 & $0{,}494505$ & $0{,}178571$ & $0{,}494505$ \\
$e_2$ & 4 & $0{,}494505$ & $-0{,}178571$ & $0{,}494505$ \\
$e_3$ & 2 & $0{,}494505$ & $-0{,}178571$ & $0{,}494505$ \\
$e_4$ & 1 & $0{,}494505$ & $0{,}178571$ & $0{,}494505$ \\
\midrule
\multicolumn{2}{l}{Suma $=\bm{K}_{55}$} & $1{,}978022$ & $0{,}000000$ & $1{,}978022$ \\
}
\expandafter\def\csname tabla:vecinos\endcsname{%
1 & $e_1$ & $-0{,}247253$ & $-0{,}178571$ & $-0{,}247253$ & $0$ & $1$ & $-0{,}178571$ & $-0{,}247253$ \\
2 & $e_1$, $e_2$ & $0{,}109890$ & $0{,}000000$ & $-0{,}604396$ & $0$ & $0$ & $0{,}000000$ & $0{,}000000$ \\
3 & $e_2$ & $-0{,}247253$ & $0{,}178571$ & $-0{,}247253$ & $0$ & $-1$ & $-0{,}178571$ & $0{,}247253$ \\
4 & $e_1$, $e_3$ & $-0{,}604396$ & $0{,}000000$ & $0{,}109890$ & $0$ & $0$ & $0{,}000000$ & $0{,}000000$ \\
6 & $e_2$, $e_4$ & $-0{,}604396$ & $0{,}000000$ & $0{,}109890$ & $0$ & $0$ & $0{,}000000$ & $0{,}000000$ \\
7 & $e_3$ & $-0{,}247253$ & $0{,}178571$ & $-0{,}247253$ & $0$ & $-1$ & $-0{,}178571$ & $0{,}247253$ \\
8 & $e_3$, $e_4$ & $0{,}109890$ & $0{,}000000$ & $-0{,}604396$ & $0$ & $0$ & $0{,}000000$ & $0{,}000000$ \\
9 & $e_4$ & $-0{,}247253$ & $-0{,}178571$ & $-0{,}247253$ & $0$ & $1$ & $-0{,}178571$ & $-0{,}247253$ \\
\midrule
\multicolumn{7}{r}{Suma $=\bm{K}_{fr}\,\bm{u}_r$} & $-0{,}714286$ & $0{,}000000$ \\
}
\expandafter\def\csname tabla:fuerzae1\endcsname{%
$g_1$ & $0{,}105662$ & $0{,}105662$ & $0{,}044658$ & $0{,}806282$ & $6{,}785722$ & $0{,}002250$ & $0{,}018940$ \\
$g_2$ & $0{,}394338$ & $0{,}105662$ & $0{,}166667$ & $2{,}339060$ & $2{,}339060$ & $0{,}024365$ & $0{,}024365$ \\
$g_3$ & $0{,}394338$ & $0{,}394338$ & $0{,}622008$ & $6{,}785722$ & $0{,}806282$ & $0{,}263799$ & $0{,}031345$ \\
$g_4$ & $0{,}105662$ & $0{,}394338$ & $0{,}166667$ & $2{,}339060$ & $2{,}339060$ & $0{,}024365$ & $0{,}024365$ \\
\midrule
\multicolumn{6}{r}{Suma en $e_1$} & $0{,}314779$ & $0{,}099015$ \\
}
\expandafter\def\csname tabla:fuerzas\endcsname{%
$e_1$ & 3 & $0{,}314779$ & $0{,}099015$ \\
$e_2$ & 4 & $0{,}314779$ & $-0{,}099015$ \\
$e_3$ & 2 & $0{,}314779$ & $-0{,}099015$ \\
$e_4$ & 1 & $0{,}314779$ & $0{,}099015$ \\
\midrule
\multicolumn{2}{l}{Suma $=(F_9,\ F_{10})$} & $1{,}259118$ & $0{,}000000$ \\
}
\expandafter\def\csname tabla:puntostres\endcsname{%
$p_{1}$ & $-0{,}774597$ & $-0{,}774597$ & $0{,}308642$ & $0{,}019290$ & $0{,}056351$ & $0{,}056351$ \\
$p_{2}$ & $-0{,}774597$ & $0{,}000000$ & $0{,}493827$ & $0{,}030864$ & $0{,}056351$ & $0{,}250000$ \\
$p_{3}$ & $-0{,}774597$ & $0{,}774597$ & $0{,}308642$ & $0{,}019290$ & $0{,}056351$ & $0{,}443649$ \\
$p_{4}$ & $0{,}000000$ & $-0{,}774597$ & $0{,}493827$ & $0{,}030864$ & $0{,}250000$ & $0{,}056351$ \\
$p_{5}$ & $0{,}000000$ & $0{,}000000$ & $0{,}790123$ & $0{,}049383$ & $0{,}250000$ & $0{,}250000$ \\
$p_{6}$ & $0{,}000000$ & $0{,}774597$ & $0{,}493827$ & $0{,}030864$ & $0{,}250000$ & $0{,}443649$ \\
$p_{7}$ & $0{,}774597$ & $-0{,}774597$ & $0{,}308642$ & $0{,}019290$ & $0{,}443649$ & $0{,}056351$ \\
$p_{8}$ & $0{,}774597$ & $0{,}000000$ & $0{,}493827$ & $0{,}030864$ & $0{,}443649$ & $0{,}250000$ \\
$p_{9}$ & $0{,}774597$ & $0{,}774597$ & $0{,}308642$ & $0{,}019290$ & $0{,}443649$ & $0{,}443649$ \\
}
\expandafter\def\csname tabla:L2e1\endcsname{%
$p_{1}$ & $0{,}012672$ & $0{,}031014$ & $0{,}787298$ & $0{,}968986$ & $0{,}033347$ & $0{,}000643$ \\
$p_{2}$ & $0{,}056219$ & $0{,}124527$ & $0{,}443649$ & $0{,}696055$ & $0{,}068375$ & $0{,}002110$ \\
$p_{3}$ & $0{,}099766$ & $0{,}173356$ & $0{,}100000$ & $0{,}173356$ & $0{,}010796$ & $0{,}000208$ \\
$p_{4}$ & $0{,}056219$ & $0{,}124527$ & $0{,}443649$ & $0{,}696055$ & $0{,}068375$ & $0{,}002110$ \\
$p_{5}$ & $0{,}249416$ & $0{,}500000$ & $0{,}250000$ & $0{,}500000$ & $0{,}125292$ & $0{,}006187$ \\
$p_{6}$ & $0{,}442613$ & $0{,}696055$ & $0{,}056351$ & $0{,}124527$ & $0{,}068881$ & $0{,}002126$ \\
$p_{7}$ & $0{,}099766$ & $0{,}173356$ & $0{,}100000$ & $0{,}173356$ & $0{,}010796$ & $0{,}000208$ \\
$p_{8}$ & $0{,}442613$ & $0{,}696055$ & $0{,}056351$ & $0{,}124527$ & $0{,}068881$ & $0{,}002126$ \\
$p_{9}$ & $0{,}785460$ & $0{,}968986$ & $0{,}012702$ & $0{,}031014$ & $0{,}034017$ & $0{,}000656$ \\
}
\expandafter\def\csname tabla:H1e1\endcsname{%
$p_{1}$ & $0{,}224877$ & $0{,}224877$ & $-1{,}774597$ & $-1{,}774597$ & $3{,}230182$ & $0{,}062311$ \\
$p_{2}$ & $0{,}997665$ & $0{,}224877$ & $-1{,}000000$ & $-1{,}774597$ & $1{,}981993$ & $0{,}061173$ \\
$p_{3}$ & $1{,}770453$ & $0{,}224877$ & $-0{,}225403$ & $-1{,}774597$ & $3{,}266734$ & $0{,}063016$ \\
$p_{4}$ & $0{,}224877$ & $0{,}997665$ & $-1{,}774597$ & $-1{,}000000$ & $1{,}981993$ & $0{,}061173$ \\
$p_{5}$ & $0{,}997665$ & $0{,}997665$ & $-1{,}000000$ & $-1{,}000000$ & $1{,}308576$ & $0{,}064621$ \\
$p_{6}$ & $1{,}770453$ & $0{,}997665$ & $-0{,}225403$ & $-1{,}000000$ & $1{,}976866$ & $0{,}061014$ \\
$p_{7}$ & $0{,}224877$ & $1{,}770453$ & $-1{,}774597$ & $-0{,}225403$ & $3{,}266734$ & $0{,}063016$ \\
$p_{8}$ & $0{,}997665$ & $1{,}770453$ & $-1{,}000000$ & $-0{,}225403$ & $1{,}976866$ & $0{,}061014$ \\
$p_{9}$ & $1{,}770453$ & $1{,}770453$ & $-0{,}225403$ & $-0{,}225403$ & $3{,}209157$ & $0{,}061905$ \\
}
\expandafter\def\csname tabla:erroreselem\endcsname{%
$e_1$ & $0{,}016376$ & $0{,}559242$ & $0{,}125000$ \\
$e_2$ & $0{,}016376$ & $0{,}559242$ & $0{,}125000$ \\
$e_3$ & $0{,}016376$ & $0{,}559242$ & $0{,}125000$ \\
$e_4$ & $0{,}016376$ & $0{,}559242$ & $0{,}125000$ \\
\midrule
Suma & $0{,}065503$ & $2{,}236968$ & $0{,}500000$ \\
}
\expandafter\def\csname tabla:refinamiento\endcsname{%
2 & $0{,}50000$ & 18 & $2{,}559362\times10^{-1}$ & $1{,}495650$ & 50 & $2{,}171756\times10^{-2}$ & $2{,}911437\times10^{-1}$ \\
4 & $0{,}25000$ & 50 & $7{,}240485\times10^{-2}$ & $7{,}247108\times10^{-1}$ & 162 & $2{,}775996\times10^{-3}$ & $7{,}263350\times10^{-2}$ \\
8 & $0{,}12500$ & 162 & $1{,}865511\times10^{-2}$ & $3{,}578143\times10^{-1}$ & 578 & $3{,}479331\times10^{-4}$ & $1{,}808843\times10^{-2}$ \\
16 & $0{,}06250$ & 578 & $4{,}698601\times10^{-3}$ & $1{,}782826\times10^{-1}$ & 2178 & $4{,}352022\times10^{-5}$ & $4{,}516070\times10^{-3}$ \\
32 & $0{,}03125$ & 2178 & $1{,}176827\times10^{-3}$ & $8{,}906145\times10^{-2}$ & 8450 & $5{,}441017\times10^{-6}$ & $1{,}128599\times10^{-3}$ \\
}
\expandafter\def\csname tabla:tasas\endcsname{%
2 $\to$ 4 & $3{,}535$ & $1{,}822$ & $2{,}064$ & $1{,}045$ & $7{,}823$ & $2{,}968$ & $4{,}008$ & $2{,}003$ \\
4 $\to$ 8 & $3{,}881$ & $1{,}957$ & $2{,}025$ & $1{,}018$ & $7{,}979$ & $2{,}996$ & $4{,}015$ & $2{,}006$ \\
8 $\to$ 16 & $3{,}970$ & $1{,}989$ & $2{,}007$ & $1{,}005$ & $7{,}995$ & $2{,}999$ & $4{,}005$ & $2{,}002$ \\
16 $\to$ 32 & $3{,}993$ & $1{,}997$ & $2{,}002$ & $1{,}001$ & $7{,}999$ & $3{,}000$ & $4{,}001$ & $2{,}001$ \\
}

% Generado por experimento_sabotaje.py: no editar a mano.
\expandafter\def\csname dato:pasaAorig\endcsname{8/8}
\expandafter\def\csname dato:pasaAsa\endcsname{8/8}
\expandafter\def\csname dato:pasaAsb\endcsname{8/8}
\expandafter\def\csname dato:pasaAsc\endcsname{0/8}
\expandafter\def\csname dato:pasaBorig\endcsname{4/4}
\expandafter\def\csname dato:pasaBsa\endcsname{2/4}
\expandafter\def\csname dato:pasaBsb\endcsname{0/4}
\expandafter\def\csname dato:vmsa\endcsname{5{,}8}
\expandafter\def\csname dato:cooksb\endcsname{18{,}17}
\expandafter\def\csname dato:cookorig\endcsname{23{,}95}
\expandafter\def\csname tabla:sabotajeA\endcsname{%
Ninguno (código original) & $2{,}00$ & $1{,}00$ & $3{,}00$ & $2{,}00$ & 8/8 \\
Rama DP devuelve la $\bm{D}$ de TP & $2{,}00$ & $1{,}00$ & $3{,}00$ & $2{,}00$ & 8/8 \\
$D_{33}$ duplicado & $2{,}00$ & $1{,}00$ & $3{,}00$ & $2{,}00$ & 8/8 \\
$D_{11}$ y $D_{22}$ +10\,\% (control) & $0{,}09$ & $0{,}27$--$0{,}31$ & $0{,}00$ & $0{,}00$ & 0/8 \\
}
\expandafter\def\csname tabla:sabotajeB\endcsname{%
Ninguno (código original) & $1{,}99$ / $1{,}00$ \cmark & $3{,}00$ / $2{,}00$ \cmark & $1{,}99$ / $1{,}00$ \cmark & $3{,}00$ / $2{,}00$ \cmark \\
Rama DP devuelve la $\bm{D}$ de TP & $1{,}99$ / $1{,}00$ \cmark & $3{,}00$ / $2{,}00$ \cmark & $-0{,}01$ / $0{,}12$ \xmark & $0{,}00$ / $0{,}00$ \xmark \\
$D_{33}$ duplicado & $-0{,}02$ / $0{,}09$ \xmark & $0{,}00$ / $0{,}00$ \xmark & $-0{,}02$ / $0{,}09$ \xmark & $0{,}00$ / $0{,}00$ \xmark \\
}


\newcommand{\cod}[1]{\path{#1}}  % monoespaciado que puede cortarse en / y .
\newcommand{\uM}{\bm{u}_M}
\newcommand{\eh}{\bm{e}}
\newcommand{\dd}{\,\mathrm{d}}

% --- Cajas --------------------------------------------------------------------------
\newtcolorbox{resultado}[1]{enhanced, breakable, colback=fondo, colframe=serieA,
  boxrule=0pt, leftrule=2.5pt, arc=0pt, outer arc=0pt, left=9pt, right=9pt,
  top=6pt, bottom=6pt, before upper={\textbf{#1}\par\smallskip}}
\newtcolorbox{aviso}[1]{enhanced, breakable, colback=fondo, colframe=apagado,
  boxrule=0pt, leftrule=2.5pt, arc=0pt, outer arc=0pt, left=9pt, right=9pt,
  top=6pt, bottom=6pt, before upper={\textbf{#1}\par\smallskip}}

% --- Gráficos -----------------------------------------------------------------------
\pgfplotsset{
  estilo/.style={
    axis line style={ejeC}, tick style={ejeC},
    every tick label/.append style={font=\footnotesize, text=tintasec},
    label style={font=\small, text=tinta}, title style={font=\small, text=tinta},
    grid=major, major grid style={grilla, line width=0.3pt},
    /pgf/number format/use comma,
  },
  colormap={secuencial}{rgb255=(205,226,251) rgb255=(158,197,244) rgb255=(109,167,236)
    rgb255=(57,135,229) rgb255=(37,106,191) rgb255=(24,79,149) rgb255=(13,54,107)},
  colormap={divergente}{rgb255=(24,79,149) rgb255=(57,135,229) rgb255=(158,197,244)
    rgb255=(240,239,236) rgb255=(241,174,168) rgb255=(215,88,83) rgb255=(137,43,42)},
}

\lstset{language=Python, basicstyle=\ttfamily\scriptsize, keywordstyle=\color{serieA}\bfseries,
  commentstyle=\color{tintasec}\itshape, stringstyle=\color{tintasec},
  numbers=left, numberstyle=\tiny\color{apagado}, numbersep=6pt, frame=single,
  rulecolor=\color{grilla}, breaklines=true, columns=fullflexible, keepspaces=true,
  showstringspaces=false, xleftmargin=1.6em, framexleftmargin=1.2em}

\begin{document}

% ==================================================================================
\begin{center}
{\LARGE\bfseries El método de las soluciones manufacturadas,\\[2pt] paso a paso}\\[8pt]
{\large Un ejemplo completo con elementos Q4 en elasticidad plana}\\[4pt]
{\small\color{tintasec} Verificación del motor de EduFEM \,·\, material de estudio \,·\, 23 de septiembre de 2026}
\end{center}
\vspace{0.4em}

\begin{aviso}{Cómo leer este documento}
\begin{itemize}[leftmargin=1.2em]
  \item El ejemplo sigue el procedimiento de Salari y Knupp~\cite[pp.~31--35]{salari} y lo
    aplica al caso base de la tesis: cuadrado unitario, tensión plana, $E=\dato{E}$,
    $\nu=\dato{nu}$ y espesor $t=\dato{t}$.
  \item La sección~\ref{sec:mano} resuelve a mano la malla más gruesa ($N=2$). Todas sus cifras
    se pueden rehacer con una calculadora; se muestran redondeadas a seis decimales, pero las
    operaciones se hicieron con todos los decimales, así que la última cifra de una suma puede
    diferir en una unidad.
  \item Las cifras las calculó \cod{generar_datos.py} con una réplica independiente del método
    (anexo~\ref{anx:replica}) que no importa nada de EduFEM. El mismo guion comprueba que
    coinciden con las de EduFEM: la matriz y el vector reducidos, el desplazamiento del nodo
    libre, las normas de $N=2$ y los 20 errores de la tabla de convergencia de la tesis.
  \item Convenciones: coma decimal; coordenadas $(x;\,y)$; grados de libertad (GDL) numerados
    desde 1, de modo que el nodo $n$ tiene los GDL $2n-1$ (dirección $x$) y $2n$ (dirección $y$).
    Las unidades son arbitrarias y consistentes: el método verifica un orden, no una magnitud.
\end{itemize}
\end{aviso}

\tableofcontents

% ==================================================================================
\section{La idea del método}\label{sec:idea}

\emph{Verificar} un código es comprobar que resuelve correctamente las ecuaciones que dice
resolver; \emph{validar} es comprobar que esas ecuaciones representan bien el fenómeno real.
Salari y Knupp adoptan la definición de Roache: un código está verificado cuando se demuestra que
su solución discreta converge a la de las ecuaciones continuas al refinar la malla, «usualmente con
cierto orden de exactitud»~\cite[p.~11]{salari}. Este documento trata solo de verificación.

Para medir el error de un código hace falta conocer la solución exacta, y los problemas de
elasticidad con solución analítica son pocos y muy particulares. El método de las soluciones
manufacturadas (MMS, por su sigla en inglés) invierte el orden del trabajo:

\begin{enumerate}[leftmargin=1.6em]
  \item se elige primero la solución, $\uM(x,y)$, una función suave cualquiera;
  \item se la sustituye en las ecuaciones de gobierno y se deduce qué datos la producen: la fuerza
    de cuerpo $\bm b$ y los valores en el contorno;
  \item se le entregan esos datos al código y se compara su resultado $\bm u_h$ con $\uM$.
\end{enumerate}

La solución manufacturada no necesita tener sentido físico: basta con que satisfaga las
ecuaciones, y eso se garantiza por construcción.

Lo que se compara no es el tamaño del error sino su ritmo de caída. En un método consistente, el
error se comporta como
\begin{equation}\label{eq:orden}
  e(h) = C\,h^{p} + \text{términos de orden superior}
\end{equation}
\cite[p.~29]{salari}, donde $h$ es el tamaño de los elementos, $C$ una constante que depende del
problema y $p$ el orden del método. La constante no dice nada sobre la corrección del código
---para la verificación, «la magnitud real del error es irrelevante»~\cite[p.~28]{salari}---;
el orden sí, porque la teoría lo fija de antemano y casi cualquier error de programación lo
altera. Por eso la prueba se repite en mallas cada vez más finas y lo que se mide es la pendiente
de $e(h)$. La figura~\ref{fig:flujo} resume los seis pasos que siguen.

\begin{figure}[H]
\centering
\begin{tikzpicture}[font=\small,
  paso/.style={draw=ejeC, fill=fondo, rounded corners=2pt, text width=3.95cm,
    minimum height=1.35cm, align=left, inner sep=5pt},
  flecha/.style={-{Stealth[length=5pt]}, tintasec, line width=0.6pt}]
  \node[paso] (p1) {\textbf{Paso 1} (§\ref{sec:paso1})\\ Ecuaciones del código y orden teórico $p$};
  \node[paso, right=0.55cm of p1] (p2) {\textbf{Paso 2} (§\ref{sec:paso2})\\ Diseño de la prueba: dominio, material, restricciones};
  \node[paso, right=0.55cm of p2] (p3) {\textbf{Paso 3} (§\ref{sec:paso3})\\ Elegir $\uM$ y deducir $\bm b$ y los valores de contorno};
  \node[paso, below=0.75cm of p3] (p4) {\textbf{Paso 4} (§\ref{sec:mano})\\ Correr el código y medir el error $e(h)$};
  \node[paso, left=0.55cm of p4] (p5) {\textbf{Paso 5} (§\ref{sec:paso5})\\ Refinar la malla: $h$, $h/2$, $h/4$, \dots};
  \node[paso, left=0.55cm of p5] (p6) {\textbf{Paso 6} (§\ref{sec:paso6})\\ Comparar el orden observado con el teórico};
  \draw[flecha] (p1) -- (p2);
  \draw[flecha] (p2) -- (p3);
  \draw[flecha] (p3) -- (p4);
  \draw[flecha] (p4) -- (p5);
  \draw[flecha] (p5) -- (p6);
\end{tikzpicture}
\caption{Los pasos del método, según Salari y Knupp~\cite[pp.~31--34]{salari}. Si el orden
observado coincide con el teórico, la prueba pasa; si no, primero se revisa la propia prueba y
después se busca el error en el código (sección~\ref{sec:resumen}).}
\label{fig:flujo}
\end{figure}

% ==================================================================================
\section{Paso 1. Las ecuaciones del código y su orden teórico}\label{sec:paso1}

«Para empezar, hay que determinar el sistema de ecuaciones de gobierno que resuelve el
código»~\cite[p.~31]{salari}. EduFEM resuelve elasticidad lineal plana con pequeñas deformaciones:

\begin{align}
  &\text{Equilibrio en } \Omega: &&
  \frac{\partial \sigma_x}{\partial x} + \frac{\partial \tau_{xy}}{\partial y} + b_x = 0, \qquad
  \frac{\partial \tau_{xy}}{\partial x} + \frac{\partial \sigma_y}{\partial y} + b_y = 0,
  \label{eq:equilibrio}\\[4pt]
  &\text{Cinemática:} &&
  \varepsilon_x = \frac{\partial u}{\partial x}, \qquad
  \varepsilon_y = \frac{\partial v}{\partial y}, \qquad
  \gamma_{xy} = \frac{\partial u}{\partial y} + \frac{\partial v}{\partial x},
  \label{eq:cinematica}\\[4pt]
  &\text{Ley constitutiva:} &&
  \begin{Bmatrix}\sigma_x\\ \sigma_y\\ \tau_{xy}\end{Bmatrix}
  = \bm D \begin{Bmatrix}\varepsilon_x\\ \varepsilon_y\\ \gamma_{xy}\end{Bmatrix},
  \label{eq:constitutiva}
\end{align}
con la matriz de tensión plana
\begin{equation}\label{eq:Dnum}
  \bm D = \frac{E}{1-\nu^2}\begin{bmatrix} 1 & \nu & 0\\ \nu & 1 & 0\\ 0 & 0 & \frac{1-\nu}{2}\end{bmatrix}
  = \begin{bmatrix} \dato{D11} & \dato{D12} & 0\\ \dato{D12} & \dato{D11} & 0\\ 0 & 0 & \dato{D33}\end{bmatrix},
\end{equation}
más las restricciones $u=\bar u$, $v=\bar v$ en la parte $\Gamma_u$ del contorno.
Conviene retener dos combinaciones de $\bm D$ que reaparecen en la sección~\ref{sec:limites}:
$D_{11}-D_{12} = E/(1+\nu) = 2\mu = \dato{DdifD}$ y $D_{33} = \mu = \dato{D33}$, donde $\mu$ es el
módulo de corte.

El método de elementos finitos (MEF) aproxima el desplazamiento dentro de cada elemento como
$\bm u_h = \sum_i N_i\,\bm u_i$, con funciones de forma $N_i$ y valores nodales $\bm u_i$, y convierte
las ecuaciones en el sistema lineal
\begin{equation}\label{eq:sistema}
  \bm K\,\bm u = \bm F, \qquad
  \bm K = \sum_e \int_{\Omega_e} \bm B^{\mathsf T}\bm D\,\bm B\; t \dd A, \qquad
  \bm F = \sum_e \int_{\Omega_e} \bm N^{\mathsf T}\bm b\; t \dd A,
\end{equation}
donde $\bm B$ contiene las derivadas de las funciones de forma ($\bm\varepsilon = \bm B\,\bm u_e$).
EduFEM evalúa esas integrales con cuadratura de Gauss de $2\times2$ puntos en el Q4 y de $3\times3$
en el Q9, los órdenes que integran exactamente la rigidez de los elementos sin
distorsionar~\cite[pp.~221--222]{hughes}.

El segundo dato del paso~1 es el orden teórico. Si las funciones de forma contienen polinomios
completos de grado $k$ y la solución es suave, el error en la norma que involucra $m$ derivadas
decrece como $h^{k+1-m}$~\cite[pp.~189--190]{hughes}\cite[pp.~106--107]{strang}. Con $m=0$ se
obtiene la norma $L^2$ del desplazamiento, y con $m=1$ la seminorma $H^1$, que mide el error de las
derivadas, es decir, de las deformaciones:

\begin{table}[H]
\centering\small
\caption{Órdenes teóricos de convergencia.}\label{tab:ordenes}
\begin{tabular}{@{}lccc@{}}
\toprule
Elemento & Grado completo $k$ & Norma $L^2$ ($h^{k+1}$) & Seminorma $H^1$ ($h^{k}$)\\
\midrule
Q4 (bilineal, 4 nodos)      & 1 & 2 & 1\\
Q9 (bicuadrático, 9 nodos)  & 2 & 3 & 2\\
\bottomrule
\end{tabular}
\end{table}

En términos prácticos: al dividir $h$ por dos, el error $L^2$ del Q4 debe dividirse por
$2^2=4$ y el del Q9 por $2^3=8$.

% ==================================================================================
\section{Paso 2. El diseño de la prueba}\label{sec:paso2}

Salari y Knupp llaman \emph{prueba de cobertura} a cada combinación de opciones del código que se
decide ejercitar~\cite[p.~31]{salari}. El guion \cod{tests/vv_mms.py} de EduFEM corre cuatro: el
cuadrado unitario y un cuadrilátero distorsionado, cada uno en tensión plana y en deformación plana.
Este documento desarrolla la configuración base:

\begin{itemize}[leftmargin=1.2em]
  \item \textbf{Dominio}: el cuadrado $\Omega=[0,1]^2$, con mallas estructuradas de $N\times N$
    elementos cuadrados de lado $h = 1/N$.
  \item \textbf{Material}: $E=\dato{E}$, $\nu=\dato{nu}$, $t=\dato{t}$, en tensión plana. Son
    valores elegidos para que todas las cifras sean de orden uno.
  \item \textbf{Restricciones}: desplazamiento prescrito en \emph{todo} el contorno. Es la
    recomendación de Salari y Knupp para la primera prueba: si los valores del contorno salen de la
    solución exacta, el error en el contorno es nulo y el que quede proviene solo de las ecuaciones
    del interior~\cite[p.~32]{salari}.
  \item \textbf{Cargas}: solo la fuerza de cuerpo $\bm b$; ni cargas nodales ni superficiales.
  \item \textbf{Elementos}: Q4, que la sección~\ref{sec:mano} resuelve a mano, y Q9, que aparece en
    el refinamiento.
\end{itemize}

% ==================================================================================
\section{Paso 3. La solución manufacturada}\label{sec:paso3}

\subsection{La elección}

La solución manufacturada de la tesis, programada en \cod{tests/vv_mms.py}, es
\begin{equation}\label{eq:uM}
  u_M(x,y) = \sin\pi x\,\sin\pi y, \qquad v_M(x,y) = \cos\pi x\,\cos\pi y .
\end{equation}
La figura~\ref{fig:campos} muestra las dos componentes. Salari y Knupp proponen cuatro pautas para
construir una solución manufacturada~\cite[p.~23]{salari}; tres se cumplen a la vista:

\begin{itemize}[leftmargin=1.2em]
  \item \textbf{Funciones analíticas suaves.} Senos y cosenos se evalúan exactamente en cualquier
    punto, y la suavidad garantiza que el orden teórico sea alcanzable.
  \item \textbf{Suficientes derivadas no triviales.} Si $\uM$ fuera lineal, el Q4 la reproduciría
    exactamente (es la prueba de la parcela), el error sería nulo en todas las mallas y no habría
    orden que medir. Una función trigonométrica no pertenece a ningún espacio polinómico, así que
    siempre deja un error de discretización que decrece con $h$.
  \item \textbf{Derivadas acotadas por una constante pequeña.} Las derivadas de $\uM$ no superan
    $\pi^2$ en valor absoluto, así que la solución no varía bruscamente dentro de un elemento.
\end{itemize}

Además, $v_M$ vale $\pm1$ en las esquinas, de modo que la prueba ejercita restricciones
\emph{no homogéneas} (desplazamientos prescritos distintos de cero). La pauta restante pide que la
solución sea «lo bastante general como para ejercitar cada término de la ecuación de
gobierno»~\cite[p.~23]{salari}; se revisa en la sección~\ref{sec:limites}, porque esta elección no
la cumple del todo.

\begin{figure}[H]
\centering
\begin{tikzpicture}
\begin{groupplot}[group style={group size=2 by 1, horizontal sep=3.1cm},
  width=0.40\textwidth, height=0.40\textwidth, view={0}{90}, estilo, grid=none,
  xmin=0, xmax=1, ymin=0, ymax=1, xtick={0,0.5,1}, ytick={0,0.5,1},
  xlabel={$x$}, ylabel={$y$}, axis on top, enlargelimits=false, colorbar,
  colorbar style={width=0.22cm, axis line style={draw=none}, ytick style={draw=none},
    yticklabel style={font=\footnotesize, text=tintasec, /pgf/number format/use comma}}]
\nextgroupplot[title={$u_M=\sin\pi x\,\sin\pi y$}, colormap name=secuencial,
  point meta min=0, point meta max=1, colorbar style={ytick={0,0.5,1}}]
\addplot3[surf, shader=interp, samples=41, samples y=41, domain=0:1, y domain=0:1]
  {sin(180*x)*sin(180*y)};
\draw[tintasec, line width=0.4pt, opacity=0.7]
  (axis cs:0.5,0,1) -- (axis cs:0.5,1,1) (axis cs:0,0.5,1) -- (axis cs:1,0.5,1);
\nextgroupplot[title={$v_M=\cos\pi x\,\cos\pi y$}, colormap name=divergente,
  point meta min=-1, point meta max=1, colorbar style={ytick={-1,0,1}}]
\addplot3[surf, shader=interp, samples=41, samples y=41, domain=0:1, y domain=0:1]
  {cos(180*x)*cos(180*y)};
\draw[tintasec, line width=0.4pt, opacity=0.7]
  (axis cs:0.5,0,1) -- (axis cs:0.5,1,1) (axis cs:0,0.5,1) -- (axis cs:1,0.5,1);
\end{groupplot}
\end{tikzpicture}
\caption{Las dos componentes de la solución manufacturada. $u_M$ se anula en todo el contorno y
alcanza su máximo, 1, en el centro; $v_M$ vale $+1$ en las esquinas $(0;0)$ y $(1;1)$, $-1$ en
$(1;0)$ y $(0;1)$, y se anula en el centro. Las líneas finas marcan la malla $N=2$ de la
sección~\ref{sec:mano}.}
\label{fig:campos}
\end{figure}

\subsection{Derivadas}

\begin{table}[H]
\centering\small
\caption{Derivadas de la solución manufacturada.}\label{tab:derivadas}
\renewcommand{\arraystretch}{1.25}
\begin{tabular}{@{}lll@{}}
\toprule
 & Componente $u_M=\sin\pi x\sin\pi y$ & Componente $v_M=\cos\pi x\cos\pi y$\\
\midrule
$\partial/\partial x$             & $\pi\cos\pi x\,\sin\pi y$     & $-\pi\sin\pi x\,\cos\pi y$\\
$\partial/\partial y$             & $\pi\sin\pi x\,\cos\pi y$     & $-\pi\cos\pi x\,\sin\pi y$\\
$\partial^2/\partial x^2$         & $-\pi^2\sin\pi x\,\sin\pi y$  & $-\pi^2\cos\pi x\,\cos\pi y$\\
$\partial^2/\partial y^2$         & $-\pi^2\sin\pi x\,\sin\pi y$  & $-\pi^2\cos\pi x\,\cos\pi y$\\
$\partial^2/\partial x\partial y$ & $\pi^2\cos\pi x\,\cos\pi y$   & $\pi^2\sin\pi x\,\sin\pi y$\\
\bottomrule
\end{tabular}
\end{table}

\subsection{Deformaciones}

Con la cinemática~\eqref{eq:cinematica} y las primeras derivadas de la tabla~\ref{tab:derivadas}:
\begin{equation}\label{eq:deformaciones}
\begin{gathered}
  \varepsilon_x = \pi\cos\pi x\,\sin\pi y, \qquad
  \varepsilon_y = -\pi\cos\pi x\,\sin\pi y,\\
  \gamma_{xy} = \pi\sin\pi x\,\cos\pi y - \pi\sin\pi x\,\cos\pi y = 0 .
\end{gathered}
\end{equation}
Aparecen dos hechos que serán importantes en la sección~\ref{sec:limites}: $\varepsilon_x +
\varepsilon_y = 0$ en todo punto (el material no cambia de área) y $\gamma_{xy} = 0$ en todo punto
(no hay distorsión angular en los ejes $x$, $y$).

\subsection{Tensiones}

Con la ley constitutiva~\eqref{eq:constitutiva}, y como $\varepsilon_y = -\varepsilon_x$:
\begin{equation}\label{eq:tensiones}
\begin{aligned}
  \sigma_x &= D_{11}\varepsilon_x + D_{12}\varepsilon_y = (D_{11}-D_{12})\,\varepsilon_x
            = 2\mu\,\pi\cos\pi x\,\sin\pi y = \dato{dosmupi}\,\cos\pi x\,\sin\pi y,\\
  \sigma_y &= D_{12}\varepsilon_x + D_{11}\varepsilon_y = -(D_{11}-D_{12})\,\varepsilon_x = -\sigma_x,\\
  \tau_{xy} &= D_{33}\,\gamma_{xy} = 0 .
\end{aligned}
\end{equation}

\subsection{La fuerza de cuerpo}

El equilibrio~\eqref{eq:equilibrio} dice qué fuerza de cuerpo hace falta para que esas tensiones
estén en equilibrio: $b_x = -(\partial\sigma_x/\partial x + \partial\tau_{xy}/\partial y)$ y
$b_y = -(\partial\tau_{xy}/\partial x + \partial\sigma_y/\partial y)$. Derivando~\eqref{eq:tensiones}:
\begin{align*}
  \frac{\partial\sigma_x}{\partial x} &= 2\mu\pi\,(-\pi\sin\pi x\,\sin\pi y) = -2\mu\pi^2\sin\pi x\,\sin\pi y
  &&\Rightarrow\quad b_x = 2\mu\pi^2\sin\pi x\,\sin\pi y,\\
  \frac{\partial\sigma_y}{\partial y} &= -2\mu\pi\,(\pi\cos\pi x\,\cos\pi y) = -2\mu\pi^2\cos\pi x\,\cos\pi y
  &&\Rightarrow\quad b_y = 2\mu\pi^2\cos\pi x\,\cos\pi y .
\end{align*}
El resultado es notablemente simple: la fuerza de cuerpo es proporcional a la propia solución,
\begin{equation}\label{eq:b}
  \bm b(x,y) = 2\mu\pi^2\,\uM(x,y) = \dato{K2}\,\bigl(\sin\pi x\,\sin\pi y;\ \cos\pi x\,\cos\pi y\bigr).
\end{equation}
Por ejemplo, en $(0,25;\,0,25)$ se tiene $u_M=v_M=\tfrac12$ y $\bm b = (\dato{bmuestra};\,\dato{bmuestra})$.

\begin{aviso}{Otra forma de llegar: la que programa EduFEM}
El guion \cod{tests/vv_mms.py} (función \cod{make_body_force_fn}) usa las ecuaciones de Navier, que
resultan de sustituir la cinemática y la ley constitutiva en el equilibrio:
\[
  b_x = -\Bigl[(\lambda+2\mu)\,u_{,xx} + \mu\,u_{,yy} + (\lambda+\mu)\,v_{,xy}\Bigr],\qquad
  b_y = -\Bigl[(\lambda+\mu)\,u_{,xy} + \mu\,v_{,xx} + (\lambda+2\mu)\,v_{,yy}\Bigr],
\]
con $\lambda = E\nu/(1-\nu^2) = \dato{lam}$ en tensión plana. Al sustituir las segundas derivadas de
la tabla~\ref{tab:derivadas} en $b_x$:
\[
  b_x = \pi^2\sin\pi x\,\sin\pi y\,\bigl[(\lambda+2\mu) + \mu - (\lambda+\mu)\bigr] = 2\mu\pi^2\sin\pi x\,\sin\pi y .
\]
El resultado es el mismo de~\eqref{eq:b}, y \emph{$\lambda$ se cancela}. Lo mismo ocurre con $b_y$.
\end{aviso}

\subsection{Los valores del contorno}

En un contorno con restricciones, los valores prescritos salen de evaluar la solución manufacturada
en los nodos correspondientes~\cite[p.~26]{salari}. En los cuatro lados del cuadrado:
\[
  \begin{array}{lll}
    y=0: & \bar u = 0, & \bar v = \cos\pi x,\\
    y=1: & \bar u = 0, & \bar v = -\cos\pi x,
  \end{array}\qquad
  \begin{array}{lll}
    x=0: & \bar u = 0, & \bar v = \cos\pi y,\\
    x=1: & \bar u = 0, & \bar v = -\cos\pi y.
  \end{array}
\]

\begin{resultado}{Resultado del paso 3: los únicos datos que recibe el código}
La fuerza de cuerpo $\bm b(x,y)$ de~\eqref{eq:b} y los valores $\bar u$, $\bar v$ en los nodos del
contorno. El código no conoce $\uM$: si la reproduce, es porque resolvió bien las ecuaciones.
\end{resultado}

% ==================================================================================
\section{Paso 4. La prueba a mano con la malla más gruesa}\label{sec:mano}

Con $N=2$ la malla tiene 9 nodos, y 8 de ellos están en el contorno con valores conocidos. Solo
queda libre el nodo central, así que el sistema por resolver es de $2\times2$. Esta sección hace a
mano todo lo que hace el código: rigidez, ensamblaje, fuerzas, restricciones, solución y medida
del error.

\subsection{Malla, numeración y valores prescritos}

\begin{figure}[H]
\centering
\begin{minipage}[b]{0.47\textwidth}\centering
\begin{tikzpicture}[x=5cm, y=5cm, font=\small]
  \foreach \cx/\cy/\e in {0.25/0.25/1, 0.75/0.25/2, 0.25/0.75/3, 0.75/0.75/4}
    \node[text=apagado] at (\cx,\cy) {$e_{\e}$};
  \draw[tintasec, line width=0.7pt] (0,0) rectangle (1,1);
  \draw[tintasec, line width=0.45pt] (0.5,0) -- (0.5,1) (0,0.5) -- (1,0.5);
  \foreach \x/\y in {0/0, 0.5/0, 1/0, 0/0.5, 1/0.5, 0/1, 0.5/1, 1/1}
    \fill[tinta] (\x,\y) circle (2.3pt);
  \filldraw[fill=white, draw=serieA, line width=1.3pt] (0.5,0.5) circle (3.3pt);
  \foreach \n/\x/\y/\anc in {1/0/0/south west, 2/0.5/0/south west, 3/1/0/south east,
      4/0/0.5/south west, 5/0.5/0.5/south west, 6/1/0.5/south east,
      7/0/1/north west, 8/0.5/1/north west, 9/1/1/north east}
    \node[anchor=\anc, font=\footnotesize, text=tinta, inner sep=3.5pt] at (\x,\y) {\n};
  \foreach \x/\l in {0/0, 0.5/{0,5}, 1/1}
    \node[below=4pt, font=\footnotesize, text=tintasec] at (\x,0) {\l};
  \foreach \y/\l in {0.5/{0,5}, 1/1}
    \node[left=4pt, font=\footnotesize, text=tintasec] at (0,\y) {\l};
  \node[font=\footnotesize, text=tintasec] at (0.5,-0.17) {$x$};
  \node[font=\footnotesize, text=tintasec] at (-0.2,0.5) {$y$};
\end{tikzpicture}\\[2pt]
{\footnotesize (a) Malla $N=2$: \tikz\fill[tinta] (0,0) circle (2.3pt); nodo del contorno,
\tikz\filldraw[fill=white, draw=serieA, line width=1.1pt] (0,0) circle (2.8pt); nodo libre.}
\end{minipage}\hfill
\begin{minipage}[b]{0.5\textwidth}\centering
\begin{tikzpicture}[x=2cm, y=2cm, font=\small]
  \draw[tintasec, line width=0.7pt] (-1,-1) rectangle (1,1);
  \draw[apagado, -{Stealth[length=4pt]}] (-1.3,0) -- (1.35,0) node[right, text=tintasec] {$\xi$};
  \draw[apagado, -{Stealth[length=4pt]}] (0,-1.3) -- (0,1.35) node[above, text=tintasec] {$\eta$};
  \foreach \a/\i in {-0.774597/0, 0/1, 0.774597/2}
    \foreach \b/\j in {-0.774597/1, 0/2, 0.774597/3}
      {\pgfmathtruncatemacro{\k}{3*\i+\j}
       \ifnum\i=0 \def\anc{north east}\else\def\anc{north west}\fi
       \filldraw[fill=white, draw=apagado, line width=0.6pt] (\a,\b) circle (2.3pt);
       \node[anchor=\anc, font=\scriptsize, text=apagado, inner sep=1.2pt] at (\a,\b) {$p_{\k}$};}
  \foreach \a/\b/\g in {-0.57735/-0.57735/1, 0.57735/-0.57735/2, 0.57735/0.57735/3, -0.57735/0.57735/4}
    {\fill[tinta] (\a,\b) circle (1.9pt);
     \node[anchor=west, font=\scriptsize, text=tinta, inner sep=2.5pt] at (\a,\b) {$g_{\g}$};}
  \foreach \x/\y/\l/\anc in {-1/-1/{1 (nodo 1)}/north east, 1/-1/{2 (nodo 2)}/north west,
      -1/1/{4 (nodo 4)}/south east}
    {\fill[tinta] (\x,\y) circle (2.3pt);
     \node[anchor=\anc, font=\footnotesize, inner sep=2pt] at (\x,\y) {\l};}
  \filldraw[fill=white, draw=serieA, line width=1.3pt] (1,1) circle (3.3pt);
  \node[anchor=south west, font=\footnotesize, inner sep=2pt] at (1,1) {3 (nodo 5)};
\end{tikzpicture}\\[2pt]
{\footnotesize (b) Elemento $e_1$ en coordenadas naturales: nodos locales (globales),
\tikz\fill[tinta] (0,0) circle (1.9pt); Gauss $2\times2$ para $\bm K$ y $\bm F$,
\tikz\filldraw[fill=white, draw=apagado, line width=0.6pt] (0,0) circle (2.3pt); Gauss $3\times3$
para el error.}
\end{minipage}
\caption{La malla del ejemplo y el elemento $e_1$ con sus puntos de integración.}
\label{fig:malla}
\end{figure}

La figura~\ref{fig:malla} muestra la malla. Los nodos se numeran fila por fila desde la esquina
inferior izquierda, como en EduFEM, y cada elemento enumera sus nodos en sentido antihorario desde
su esquina inferior izquierda: $e_1=(1,2,5,4)$, $e_2=(2,3,6,5)$, $e_3=(4,5,8,7)$ y $e_4=(5,6,9,8)$.
Los GDL libres son el 9 y el 10 ($u_5$ y $v_5$); los otros 16 están restringidos. La
tabla~\ref{tab:nodos} da los valores prescritos, que salen de evaluar~\eqref{eq:uM} en cada nodo.

\begin{table}[H]
\centering\small
\caption{Nodos de la malla $N=2$ y valores prescritos $\bar u = u_M$, $\bar v = v_M$.}\label{tab:nodos}
\begin{tabular}{@{}cccccc@{}}
\toprule
Nodo & $x$ & $y$ & Tipo & $\bar u$ & $\bar v$\\
\midrule
\tabla{nodos}
\bottomrule
\end{tabular}
\end{table}

\subsection{Funciones de forma, jacobiano y matriz \texorpdfstring{$\bm B$}{B}}

Las funciones de forma del Q4 en las coordenadas naturales $\xi,\eta\in[-1,1]$ son
\begin{equation}\label{eq:formas}
  N_1 = \tfrac14(1-\xi)(1-\eta),\quad N_2 = \tfrac14(1+\xi)(1-\eta),\quad
  N_3 = \tfrac14(1+\xi)(1+\eta),\quad N_4 = \tfrac14(1-\xi)(1+\eta).
\end{equation}
En el elemento $e_1$, el mapeo $x=\sum N_i x_i$, $y=\sum N_i y_i$ se reduce a $x = 0,25\,(1+\xi)$ e
$y = 0,25\,(1+\eta)$, de modo que
\[
  \bm J = \begin{bmatrix}\partial x/\partial\xi & \partial y/\partial\xi\\
                         \partial x/\partial\eta & \partial y/\partial\eta\end{bmatrix}
        = \begin{bmatrix}\dato{J11} & 0\\ 0 & \dato{J11}\end{bmatrix},\qquad
  \det\bm J = \dato{detJ},\qquad
  \frac{\partial N_i}{\partial x} = \dato{Jinv}\,\frac{\partial N_i}{\partial \xi},\quad
  \frac{\partial N_i}{\partial y} = \dato{Jinv}\,\frac{\partial N_i}{\partial \eta}.
\]
Los otros tres elementos son traslaciones de $e_1$ y tienen el mismo jacobiano. La matriz $\bm B$
($3\times8$) ordena esas derivadas de modo que $\bm\varepsilon = \bm B\,\bm u_e$, con
$\bm u_e = (u_1,v_1,u_2,v_2,u_3,v_3,u_4,v_4)$:
\[
  \bm B = \begin{bmatrix}
    N_{1,x} & 0 & N_{2,x} & 0 & N_{3,x} & 0 & N_{4,x} & 0\\
    0 & N_{1,y} & 0 & N_{2,y} & 0 & N_{3,y} & 0 & N_{4,y}\\
    N_{1,y} & N_{1,x} & N_{2,y} & N_{2,x} & N_{3,y} & N_{3,x} & N_{4,y} & N_{4,x}
  \end{bmatrix}.
\]
Los cuatro puntos de Gauss $2\times2$ están en $\xi,\eta=\pm1/\sqrt3=\pm\dato{gdos}$, con peso $w=1$;
se numeran $g_1,\dots,g_4$ en sentido antihorario desde $(-,-)$ (figura~\ref{fig:malla}b). En
$g_1=(-\dato{gdos};\,-\dato{gdos})$, por ejemplo:

\begin{table}[H]
\centering\small
\caption{Matriz $\bm B$ del elemento $e_1$ en el punto $g_1$.}\label{tab:B}
\setlength{\tabcolsep}{5pt}
\begin{tabular}{@{}*{8}{r}@{}}
\toprule
$u_1$ & $v_1$ & $u_2$ & $v_2$ & $u_3$ & $v_3$ & $u_4$ & $v_4$\\
\midrule
\tabla{B}
\bottomrule
\end{tabular}
\end{table}

\subsection{La rigidez del elemento}

La rigidez de cada elemento es la integral de~\eqref{eq:sistema} evaluada con los cuatro puntos de
Gauss:
\begin{equation}\label{eq:ke}
  \bm k_e = \sum_{g=1}^{4} \bm B_g^{\mathsf T}\,\bm D\,\bm B_g\;\det\bm J\;w_g\;t .
\end{equation}
La tabla~\ref{tab:ke} muestra el resultado. Los cuatro elementos tienen la misma $\bm k_e$
(tienen igual forma y orientación). En un cuadrado, además, $\bm k_e$ no depende del lado $a$:
$\bm B$ es proporcional a $1/a$ y $\det\bm J$ a $a^2$, y los dos factores se compensan.

\begin{table}[H]
\centering\small
\caption{Matriz de rigidez $\bm k_e$ de cada elemento (idéntica en los cuatro).}\label{tab:ke}
\setlength{\tabcolsep}{4.2pt}
\begin{tabular}{@{}c*{8}{r}@{}}
\toprule
 & $u_1$ & $v_1$ & $u_2$ & $v_2$ & $u_3$ & $v_3$ & $u_4$ & $v_4$\\
\midrule
\tabla{ke}
\bottomrule
\end{tabular}
\end{table}

Tres comprobaciones que se pueden hacer a mano:
\begin{itemize}[leftmargin=1.2em]
  \item $\bm k_e$ es simétrica.
  \item En cada fila, la suma de las columnas $u$ y la suma de las columnas $v$ son nulas: una
    traslación rígida no genera fuerzas. En la primera fila:
    $\dato{ck1}+(\dato{ck3})+(\dato{ck5})+\dato{ck7} = 0$.
  \item Cada entrada es uno de estos ocho valores multiplicado por $E\,t/(1-\nu^2) = \dato{c}$:
    \par\smallskip
    {\small\renewcommand{\arraystretch}{1.35}
    \begin{tabular}{@{}l*{8}{c}@{}}
      \toprule
      & $k_1$ & $k_2$ & $k_3$ & $k_4$ & $k_5$ & $k_6$ & $k_7$ & $k_8$\\
      \midrule
      Fórmula & $\tfrac12-\tfrac\nu6$ & $\tfrac18+\tfrac\nu8$ & $-\tfrac14-\tfrac{\nu}{12}$ &
        $-\tfrac18+\tfrac{3\nu}8$ & $-\tfrac14+\tfrac{\nu}{12}$ & $-\tfrac18-\tfrac\nu8$ &
        $\tfrac\nu6$ & $\tfrac18-\tfrac{3\nu}8$\\
      $\nu=\dato{nu}$ & $\dato{k1}$ & $\dato{k2}$ & $\dato{k3}$ & $\dato{k4}$ & $\dato{k5}$ &
        $\dato{k6}$ & $\dato{k7}$ & $\dato{k8}$\\
      \bottomrule
    \end{tabular}}
    \par\smallskip
    La diagonal, por ejemplo, vale $\dato{c}\times\dato{k1} = \dato{ck1}$.
\end{itemize}

\subsection{Las ecuaciones del nodo libre}

Solo interesan las filas de $\bm K$ que corresponden a los GDL libres, la 9 y la 10. El nodo 5
pertenece a los cuatro elementos, con un número local distinto en cada uno (3 en $e_1$, 4 en $e_2$,
2 en $e_3$ y 1 en $e_4$), y su bloque diagonal $\bm K_{55}$ ($2\times2$) es la suma de los cuatro
bloques diagonales correspondientes:

\begin{table}[H]
\centering\small
\caption{Ensamblaje del bloque diagonal $\bm K_{55}$.}\label{tab:Kdiag}
\begin{tabular}{@{}ccrrr@{}}
\toprule
Elemento & Nodo local & $K_{xx}$ & $K_{xy}=K_{yx}$ & $K_{yy}$\\
\midrule
\tabla{Kdiag}
\bottomrule
\end{tabular}
\end{table}

Los términos cruzados $\pm\dato{ck2}$ tienen signos alternados y se cancelan, así que
$\bm K_{55} = \dato{Kcc}\,\bm I$. Las filas 9 y 10 contienen también los bloques $\bm K_{5j}$ que
acoplan el nodo 5 con cada nodo del contorno. La tabla~\ref{tab:vecinos} los reúne junto con los
valores prescritos de cada nodo, que se usarán en la sección~\ref{sec:restricciones}.

\begin{table}[H]
\centering\small
\caption{Bloques de acoplamiento $\bm K_{5j}$ entre el nodo libre y los nodos del contorno, valores
prescritos y producto $\bm K_{5j}\,\bar{\bm u}_j$. Todos los bloques son simétricos.}\label{tab:vecinos}
\setlength{\tabcolsep}{4.5pt}
\begin{tabular}{@{}clrrrrrrr@{}}
\toprule
 & & \multicolumn{3}{c}{Bloque $\bm K_{5j}$} & \multicolumn{2}{c}{Prescritos} & \multicolumn{2}{c}{$\bm K_{5j}\,\bar{\bm u}_j$}\\
\cmidrule(lr){3-5}\cmidrule(lr){6-7}\cmidrule(l){8-9}
Nodo $j$ & Elementos & $K_{xx}$ & $K_{xy}$ & $K_{yy}$ & $\bar u_j$ & $\bar v_j$ & fila 9 & fila 10\\
\midrule
\tabla{vecinos}
\bottomrule
\end{tabular}
\end{table}

Comprobación: la fila completa del nodo 5 suma cero, porque una traslación rígida de toda la
malla no genera fuerzas. En la dirección $x$:
$\dato{Kcc} + 4(\dato{ck5}) + 2(\dato{Kvecx}) + 2(\dato{Kvecy}) = 0$.

\subsection{Las fuerzas nodales de la fuerza de cuerpo}

Cada elemento aporta el vector $\bm f_e = \sum_g \bm N(\xi_g,\eta_g)^{\mathsf T}\,\bm b(x_g,y_g)\,
\det\bm J\,w_g\,t$, y del vector ensamblado solo importan las componentes del nodo libre. En $e_1$,
el nodo 5 es el nodo local 3, cuya función de forma es $N_3 = \tfrac14(1+\xi)(1+\eta)$. Los puntos
de Gauss caen en $x,y\in\{\dato{gA};\ \dato{gB}\}$ $(=0,25\mp0,25/\sqrt3)$, y $N_3$ vale
$\dato{Ncerca}$ en el punto más cercano al nodo 5 ($g_3$), $\dato{Nmixto}$ en los dos intermedios y
$\dato{Nlejos}$ en el más lejano ($g_1$):

\begin{table}[H]
\centering\small
\caption{Aporte del elemento $e_1$ a las fuerzas del nodo 5 ($\det\bm J = \dato{detJ}$, $w_g=1$).}\label{tab:fuerzae1}
\setlength{\tabcolsep}{5pt}
\begin{tabular}{@{}c*{7}{r}@{}}
\toprule
Punto & $x$ & $y$ & $N_3$ & $b_x$ & $b_y$ & $N_3 b_x \det\bm J\,w$ & $N_3 b_y\det\bm J\,w$\\
\midrule
\tabla{fuerzae1}
\bottomrule
\end{tabular}
\end{table}

Para rehacer la componente $x$ con calculadora, conviene sacar factor común $2\mu\pi^2\det\bm J$:
con $\sin(\pi\cdot\dato{gA}) = \dato{sA}$ y $\sin(\pi\cdot\dato{gB}) = \dato{sB}$,
\begin{align*}
  \sum_g N_3\,\sin\pi x_g\,\sin\pi y_g
  &= \dato{Ncerca}\cdot\dato{sB}^2 + 2\cdot\dato{Nmixto}\cdot\dato{sB}\cdot\dato{sA}\\
  &\quad + \dato{Nlejos}\cdot\dato{sA}^2 = \dato{sumaNss},\\
  f_x^{(e_1)} &= \dato{K2}\times\dato{sumaNss}\times\dato{detJ} = \dato{fe1x}.
\end{align*}
Los cuatro elementos aportan lo mismo en $x$ por simetría. En $y$, en cambio, los aportes se
cancelan de a pares, porque $\cos\pi x$ cambia de signo en $x=0,5$ y $\cos\pi y$ en $y=0,5$:

\begin{table}[H]
\centering\small
\caption{Aporte de cada elemento a las fuerzas del nodo 5.}\label{tab:fuerzas}
\begin{tabular}{@{}ccrr@{}}
\toprule
Elemento & Nodo local & Fila 9 ($x$) & Fila 10 ($y$)\\
\midrule
\tabla{fuerzas}
\bottomrule
\end{tabular}
\end{table}

\subsection{Las restricciones no homogéneas}\label{sec:restricciones}

Si se separan los GDL libres ($f$) de los restringidos ($r$), el sistema~\eqref{eq:sistema} queda
\[
  \begin{bmatrix}\bm K_{ff} & \bm K_{fr}\\ \bm K_{rf} & \bm K_{rr}\end{bmatrix}
  \begin{Bmatrix}\bm u_f\\ \bm u_r\end{Bmatrix}
  = \begin{Bmatrix}\bm F_f\\ \bm F_r + \bm R_r\end{Bmatrix},
\]
donde $\bm R_r$ son las reacciones, desconocidas. La primera fila contiene solo incógnitas en
$\bm u_f$, porque $\bm u_r$ es conocido; al pasar el término conocido a la derecha:
\begin{equation}\label{eq:reducido}
  \bm K_{ff}\,\bm u_f = \bm F_f - \bm K_{fr}\,\bm u_r .
\end{equation}
Esta es la \emph{sustitución estática} que hace el solucionador de EduFEM (función
\cod{apply_boundary_conditions} de \cod{fem/solver.py}). En el ejemplo,
\[
  \bm K_{ff} = \bm K_{55},\qquad
  \bm F_f = \begin{Bmatrix}F_9\\ F_{10}\end{Bmatrix} = \begin{Bmatrix}\dato{F9}\\ \dato{F10}\end{Bmatrix},\qquad
  \bm K_{fr}\,\bm u_r = \sum_j \bm K_{5j}\,\bar{\bm u}_j = \begin{Bmatrix}\dato{Kfrx}\\ \dato{Kfry}\end{Bmatrix},
\]
donde el último vector es la fila de sumas de la tabla~\ref{tab:vecinos}. En la fila 9 solo aportan
las esquinas: su término cruzado $\pm\dato{nesq}$, multiplicado por $\bar v_j=\pm1$, da siempre
$-\dato{nesq}$. Entonces
\[
  \bm F_f - \bm K_{fr}\,\bm u_r = \begin{Bmatrix}\dato{F9} - (\dato{Kfrx})\\ \dato{F10} - \dato{Kfry}\end{Bmatrix}
  = \begin{Bmatrix}\dato{Fredx}\\ \dato{Fredy}\end{Bmatrix}.
\]
Más de un tercio del lado derecho proviene de los desplazamientos prescritos: así se ejercitan las
restricciones no homogéneas.

\subsection{La solución y el error en el nodo}

\[
  \begin{bmatrix}\dato{Kcc} & 0\\ 0 & \dato{Kcc}\end{bmatrix}
  \begin{Bmatrix}u_5\\ v_5\end{Bmatrix} = \begin{Bmatrix}\dato{Fredx}\\ \dato{Fredy}\end{Bmatrix}
  \quad\Longrightarrow\quad
  u_5 = \frac{\dato{Fredx}}{\dato{Kcc}} = \dato{u5},\qquad v_5 = 0 .
\]
La solución exacta en el nodo es $\uM(0,5;\,0,5) = (1;\,0)$. El error nodal es
$u_5 - 1 = \dato{erru5}$, es decir, del $\dato{erru5pct}$\,\%, y $v_5$ es exacto.

Una última comprobación, de equilibrio global: las reacciones $\bm R = \bm K\bm u - \bm F$ en los GDL
restringidos deben equilibrar la fuerza de cuerpo total. La suma de las fuerzas nodales es
$\sum F_x = \dato{sumaFx}$ y $\sum F_y = \dato{sumaFy}$, y la de las reacciones es
$\sum R_x = \dato{sumaRx}$ y $\sum R_y = \dato{sumaRy}$.

\subsection{El error en todo el dominio}\label{sec:normas}

Un error nodal pequeño no significa que la solución sea buena en todas partes. Por eso el método
mide el error integrado sobre todo el dominio, con dos normas:
\begin{equation}\label{eq:normas}
  \|\eh\|_{L^2} = \Bigl(\int_\Omega \bigl[(u_h-u_M)^2 + (v_h-v_M)^2\bigr]\dd A\Bigr)^{1/2},\qquad
  |\eh|_{H^1} = \Bigl(\int_\Omega \sum_{c,j}\Bigl(\frac{\partial (u_h-u_M)_c}{\partial x_j}\Bigr)^2\dd A\Bigr)^{1/2},
\end{equation}
donde la suma recorre las cuatro derivadas ($\partial u/\partial x$, $\partial u/\partial y$,
$\partial v/\partial x$ y $\partial v/\partial y$). La primera mide el error del desplazamiento, y
la segunda el de sus derivadas, que son las que determinan deformaciones y tensiones.

Las integrales se evalúan con un punto de Gauss más por dirección que la rigidez: $3\times3$ en el
Q4. Una cuadratura ajustada a la rigidez muestrea el campo justo en los puntos donde el error de
discretización tiende a anularse, y subestimaría el error (en EduFEM:
\cod{fem/error_norms.py}). Los puntos $3\times3$ están en $\xi,\eta\in\{-\sqrt{0,6};\,0;\,
\sqrt{0,6}\}$, con pesos $\tfrac59$, $\tfrac89$ y $\tfrac59$ (tabla~\ref{tab:puntos}).

Dentro de $e_1$ la solución de elementos finitos es muy sencilla: de sus cuatro nodos (1, 2, 5, 4),
solo el 5 tiene $u\neq0$ y solo el 1 tiene $v\neq0$ ($\bar v_1 = 1$; $v_5=0$). Por lo tanto
\[
  u_h = u_5\,N_3 = \dato{u5}\cdot\tfrac14(1+\xi)(1+\eta), \qquad v_h = N_1 = \tfrac14(1-\xi)(1-\eta).
\]

\begin{table}[H]
\centering\small
\caption{Puntos de Gauss $3\times3$ en el elemento $e_1$ ($w$: peso; $w\det\bm J$: área que representa cada punto).}\label{tab:puntos}
\begin{tabular}{@{}c*{6}{r}@{}}
\toprule
Punto & $\xi$ & $\eta$ & $w$ & $w\det\bm J$ & $x$ & $y$\\
\midrule
\tabla{puntostres}
\bottomrule
\end{tabular}
\end{table}

\begin{table}[H]
\centering\small
\caption{Norma $L^2$: aporte de cada punto de $e_1$. La última columna es $[(u_h-u_M)^2+(v_h-v_M)^2]\,w\det\bm J$.}\label{tab:L2}
\begin{tabular}{@{}c*{6}{r}@{}}
\toprule
Punto & $u_h$ & $u_M$ & $v_h$ & $v_M$ & $(u_h-u_M)^2+(v_h-v_M)^2$ & Aporte\\
\midrule
\tabla{L2e1}
\bottomrule
\end{tabular}
\end{table}

Por ejemplo, en el centro del elemento ($p_5$: $\xi=\eta=0$, $x=y=\dato{cx}$): $u_h =
\dato{u5}\cdot\tfrac14 = \dato{uhc}$ frente a $u_M = \dato{uMc}$, y $v_h = \tfrac14 = \dato{vhc}$ frente a
$v_M = \dato{vMc}$; la suma de cuadrados es $\dato{d2c}$ y, multiplicada por $w\det\bm J =
\tfrac{64}{81}\cdot\dato{detJ} = \dato{dAc}$, aporta $\dato{cL2c}$.

Para la seminorma $H^1$ se derivan los campos anteriores ($\partial/\partial x = 4\,\partial/\partial\xi$):
\[
  \frac{\partial u_h}{\partial x} = u_5\,(1+\eta),\quad
  \frac{\partial u_h}{\partial y} = u_5\,(1+\xi),\quad
  \frac{\partial v_h}{\partial x} = -(1-\eta),\quad
  \frac{\partial v_h}{\partial y} = -(1-\xi),
\]
y se comparan con las derivadas exactas de la tabla~\ref{tab:derivadas}. En $p_5$, con las
filas $(\partial/\partial x,\ \partial/\partial y)$ de $u$ y de $v$:
\[
  \nabla\bm u_h = \begin{bmatrix}\dato{gux} & \dato{guy}\\ \dato{gvx} & \dato{gvy}\end{bmatrix},\qquad
  \nabla\uM = \begin{bmatrix}\pi/2 & \pi/2\\ -\pi/2 & -\pi/2\end{bmatrix}
            = \begin{bmatrix}\dato{eux} & \dato{euy}\\ \dato{evx} & \dato{evy}\end{bmatrix}.
\]
Las diferencias son $\dato{dux}$ (dos veces) y $\dato{dvx}$ (dos veces), la suma de sus cuadrados
es $\dato{e2c}$ y el aporte es $\dato{e2c}\times\dato{dAc} = \dato{cH1c}$.

\begin{table}[H]
\centering\small
\caption{Seminorma $H^1$: aporte de cada punto de $e_1$. $|\nabla\eh|^2$ es la suma de los cuadrados de las cuatro diferencias entre el gradiente calculado y el exacto.}\label{tab:H1}
\begin{tabular}{@{}c*{6}{r}@{}}
\toprule
Punto & $\partial u_h/\partial x$ & $\partial u_h/\partial y$ & $\partial v_h/\partial x$ & $\partial v_h/\partial y$ & $|\nabla\eh|^2$ & Aporte\\
\midrule
\tabla{H1e1}
\bottomrule
\end{tabular}
\end{table}

Los otros tres elementos son imágenes especulares de $e_1$ y aportan exactamente lo mismo
(tabla~\ref{tab:elementos}).

\begin{table}[H]
\centering\small
\caption{Integrales por elemento: cuadrado del error $L^2$, cuadrado del error $H^1$ y cuadrado de la norma de $\uM$.}\label{tab:elementos}
\begin{tabular}{@{}crrr@{}}
\toprule
Elemento & $\int (\bm u_h-\uM)^2\dd A$ & $\int |\nabla\eh|^2\dd A$ & $\int \uM^2\dd A$\\
\midrule
\tabla{erroreselem}
\bottomrule
\end{tabular}
\end{table}

\begin{resultado}{Resultado del paso 4 con $N=2$}
Con las sumas de la tabla~\ref{tab:elementos}:
\begin{gather*}
  \|\eh\|_{L^2} = \sqrt{\dato{totL2}} = \dato{L2}, \qquad
  |\eh|_{H^1} = \sqrt{\dato{totH1}} = \dato{H1},\\
  \frac{\|\eh\|_{L^2}}{\|\uM\|_{L^2}} = \frac{\dato{L2}}{\dato{normaEx}} = \dato{L2rel}\ \text{\%}.
\end{gather*}
El desplazamiento en el nodo es casi exacto ($\dato{erru5pct}$\,\% de error), pero el campo tiene un
error del $\dato{L2rel}$\,\%: entre nodos, el Q4 interpola con rectas una función senoidal, y con
cuatro elementos esa aproximación es tosca. Estos dos números son el primer punto de las curvas de
convergencia de la sección~\ref{sec:paso6}.
\end{resultado}

% ==================================================================================
\section{Paso 5. El refinamiento de la malla}\label{sec:paso5}

El paso 4 se repite con $N = 4$, 8, 16 y 32. A mano ya no es viable (con $N=32$ el Q4 tiene
2178~GDL), pero el procedimiento es idéntico: la réplica del anexo~\ref{anx:replica} ejecuta en un
bucle exactamente las operaciones de la sección~\ref{sec:mano}. Se agrega el Q9, cuyas funciones de
forma son bicuadráticas (nueve nodos por elemento), con rigidez integrada con Gauss $3\times3$ y
error con $4\times4$.

\begin{table}[H]
\centering\small
\caption{Errores de la solución de elementos finitos al refinar la malla ($h=1/N$).}\label{tab:refinamiento}
\setlength{\tabcolsep}{4.5pt}
\begin{tabular}{@{}cc rcc rcc@{}}
\toprule
 & & \multicolumn{3}{c}{Q4} & \multicolumn{3}{c}{Q9}\\
\cmidrule(lr){3-5}\cmidrule(l){6-8}
$N$ & $h$ & GDL & $\|\eh\|_{L^2}$ & $|\eh|_{H^1}$ & GDL & $\|\eh\|_{L^2}$ & $|\eh|_{H^1}$\\
\midrule
\tabla{refinamiento}
\bottomrule
\end{tabular}
\end{table}

Estas cifras coinciden hasta la séptima cifra significativa con las de EduFEM
(\cod{docs/vyv/datos/mms_q4.csv} y \cod{mms_q9.csv}); la tabla del MMS de la tesis muestra las
mismas con cuatro cifras.

% ==================================================================================
\section{Paso 6. El orden observado frente al teórico}\label{sec:paso6}

Si $e(h) = C\,h^p$, al pasar de una malla de tamaño $h$ a otra de tamaño $h/r$ el cociente de errores
es $e(h)/e(h/r) = r^p$, y el orden observado se despeja como~\cite[ec.~25, p.~30]{salari}
\begin{equation}\label{eq:tasa}
  p_{\text{obs}} = \frac{\ln\bigl(e_{\text{gruesa}}/e_{\text{fina}}\bigr)}{\ln r}
  \;\overset{r=2}{=}\; \log_2\frac{e_N}{e_{2N}} .
\end{equation}
Por ejemplo, para el Q4 en norma $L^2$ entre $N=16$ y $N=32$:
\[
  \frac{e_{16}}{e_{32}} = \frac{\dato{eL2dieciseis}}{\dato{eL2treintaydos}} = \dato{cocienteL2},
  \qquad p_{\text{obs}} = \log_2\dato{cocienteL2} = \frac{\ln\dato{cocienteL2}}{\ln 2} = \dato{tasaL2}.
\]

\begin{table}[H]
\centering\small
\caption{Cociente de errores entre mallas sucesivas y orden observado~\eqref{eq:tasa}.}\label{tab:tasas}
\setlength{\tabcolsep}{4pt}
\begin{tabular}{@{}l rr rr rr rr@{}}
\toprule
 & \multicolumn{2}{c}{Q4, $L^2$} & \multicolumn{2}{c}{Q4, $H^1$} & \multicolumn{2}{c}{Q9, $L^2$} & \multicolumn{2}{c}{Q9, $H^1$}\\
\cmidrule(lr){2-3}\cmidrule(lr){4-5}\cmidrule(lr){6-7}\cmidrule(l){8-9}
Mallas & cociente & orden & cociente & orden & cociente & orden & cociente & orden\\
\midrule
\tabla{tasas}
\midrule
Teórico & $4$ & $2$ & $2$ & $1$ & $8$ & $3$ & $4$ & $2$\\
\bottomrule
\end{tabular}
\end{table}

\begin{figure}[H]
\centering
\begin{tikzpicture}
\begin{groupplot}[group style={group size=2 by 1, horizontal sep=1.7cm},
  width=0.47\textwidth, height=0.44\textwidth, estilo, xmode=log, ymode=log,
  xmin=0.024, xmax=0.9,
  xtick={0.03125,0.0625,0.125,0.25,0.5},
  xticklabels={$\frac{1}{32}$,$\frac{1}{16}$,$\frac18$,$\frac14$,$\frac12$},
  xlabel={tamaño de los elementos $h$},
  legend style={draw=none, fill=white, fill opacity=0.85, text opacity=1, font=\footnotesize, text=tintasec},
  legend cell align=left]
\nextgroupplot[ylabel={$\|\eh\|_{L^2}$}, title={Desplazamiento (norma $L^2$)}, legend pos=south east]
\addplot[serieA, line width=1.1pt, mark=*, mark size=2.2pt, mark options={draw=white, line width=0.5pt, fill=serieA}]
  table[x=h, y=L2q4] {convergencia.dat};
\addlegendentry{Q4}
\addplot[serieB, line width=1.1pt, mark=square*, mark size=2.1pt, mark options={draw=white, line width=0.5pt, fill=serieB}]
  table[x=h, y=L2q9] {convergencia.dat};
\addlegendentry{Q9}
\addplot[apagado, densely dashed, line width=0.6pt, forget plot] table[x=h, y=L2q4] {pendientes.dat}
  node[pos=0.55, sloped, below, font=\scriptsize, text=tintasec] {pendiente 2};
\addplot[apagado, densely dashed, line width=0.6pt, forget plot] table[x=h, y=L2q9] {pendientes.dat}
  node[pos=0.55, sloped, below, font=\scriptsize, text=tintasec] {pendiente 3};
\node[right, font=\footnotesize, text=tintasec] at (axis cs:0.52,2.559362e-01) {Q4};
\node[right, font=\footnotesize, text=tintasec] at (axis cs:0.52,2.171756e-02) {Q9};
\nextgroupplot[ylabel={$|\eh|_{H^1}$}, title={Deformaciones (seminorma $H^1$)}, legend pos=south east]
\addplot[serieA, line width=1.1pt, mark=*, mark size=2.2pt, mark options={draw=white, line width=0.5pt, fill=serieA}]
  table[x=h, y=H1q4] {convergencia.dat};
\addlegendentry{Q4}
\addplot[serieB, line width=1.1pt, mark=square*, mark size=2.1pt, mark options={draw=white, line width=0.5pt, fill=serieB}]
  table[x=h, y=H1q9] {convergencia.dat};
\addlegendentry{Q9}
\addplot[apagado, densely dashed, line width=0.6pt, forget plot] table[x=h, y=H1q4] {pendientes.dat}
  node[pos=0.55, sloped, below, font=\scriptsize, text=tintasec] {pendiente 1};
\addplot[apagado, densely dashed, line width=0.6pt, forget plot] table[x=h, y=H1q9] {pendientes.dat}
  node[pos=0.55, sloped, below, font=\scriptsize, text=tintasec] {pendiente 2};
\node[right, font=\footnotesize, text=tintasec] at (axis cs:0.52,1.495650) {Q4};
\node[right, font=\footnotesize, text=tintasec] at (axis cs:0.52,2.911437e-01) {Q9};
\end{groupplot}
\end{tikzpicture}
\caption{Convergencia en escala logarítmica. En estos ejes, $e = C\,h^p$ es una recta de pendiente
$p$; las rectas discontinuas marcan las pendientes teóricas, desplazadas hacia abajo para no tapar
los datos.}
\label{fig:convergencia}
\end{figure}

Lectura de los resultados (tabla~\ref{tab:tasas} y figura~\ref{fig:convergencia}):

\begin{itemize}[leftmargin=1.2em]
  \item Los órdenes observados tienden a 2 y 1 en el Q4 y a 3 y 2 en el Q9, los valores teóricos de
    la tabla~\ref{tab:ordenes}. Duplicar $N$ divide el error $L^2$ del Q4 por 4 y el del Q9 por 8.
  \item En las mallas gruesas el orden todavía no se estabiliza (entre $N=2$ y $N=4$, el Q4 da
    $\dato{tasaPrimeraQcuatro}$ en $L^2$): allí los términos de orden superior de~\eqref{eq:orden}
    todavía pesan. Por eso el criterio se aplica al par más fino. Tampoco hay que esperar recuperar
    el orden teórico con más de dos o tres cifras significativas~\cite[p.~30]{salari}.
  \item El criterio de la tesis es $|p_{\text{obs}} - p_{\text{teórico}}| < 0,5$ entre $N=16$ y
    $N=32$. Ese margen separa el orden esperado del inmediato inferior, que es lo que típicamente
    produce un error de programación.
\end{itemize}

\begin{table}[H]
\centering\small
\caption{Veredicto de la prueba (par de mallas $N=16\to32$).}\label{tab:veredicto}
\begin{tabular}{@{}llccc@{}}
\toprule
Elemento & Norma & Orden observado & Orden teórico & $|p_{\text{obs}}-p|<0,5$\\
\midrule
Q4 & $L^2$ & $\dato{tasaQcuatroLdos}$ & 2 & \cmark\ pasa\\
Q4 & $H^1$ & $\dato{tasaQcuatroHuno}$ & 1 & \cmark\ pasa\\
Q9 & $L^2$ & $\dato{tasaQnueveLdos}$ & 3 & \cmark\ pasa\\
Q9 & $H^1$ & $\dato{tasaQnueveHuno}$ & 2 & \cmark\ pasa\\
\bottomrule
\end{tabular}
\end{table}

EduFEM mide además el error $L^2$ del campo de tensiones \emph{recuperado} $\bm\sigma^*$ (tensiones
en los puntos de Gauss, extrapoladas a los nodos y promediadas), que es el que muestra el
post-proceso. Salari y Knupp advierten que el orden de las magnitudes derivadas puede ser menor que
el de la solución~\cite[p.~29]{salari}. Como la teoría no fija un orden para esa cadena de
recuperación, la tesis usa una cota empírica mínima (1,4 en Q4 y 1,9 en Q9); los órdenes observados
son 1,54 y 2,00 (\cod{docs/vyv/datos/mms_resumen.csv}).

% ==================================================================================
\section{Qué verificó la prueba y qué no}\label{sec:limites}

Una prueba superada verifica las partes del código que la solución manufacturada ejercita, y
ninguna otra. Por eso la segunda pauta de Salari y Knupp pide una solución «lo bastante general como
para ejercitar cada término de la ecuación de gobierno»~\cite[p.~23]{salari}. Para comprobarlo con
la ley constitutiva, basta con repetir el cálculo de la sección~\ref{sec:paso3} con una matriz
genérica de la misma forma:
\[
  \bm D = \begin{bmatrix} D_{11} & D_{12} & 0\\ D_{12} & D_{11} & 0\\ 0 & 0 & D_{33}\end{bmatrix}.
\]
Con las deformaciones~\eqref{eq:deformaciones} ($\varepsilon_y = -\varepsilon_x$ y $\gamma_{xy}=0$)
se obtiene $\sigma_x = (D_{11}-D_{12})\,\varepsilon_x$, $\sigma_y = -\sigma_x$ y $\tau_{xy} =
D_{33}\cdot 0 = 0$, de modo que
\begin{equation}\label{eq:bgenerica}
  \bm b = \pi^2\,(D_{11}-D_{12})\,\bigl(\sin\pi x\,\sin\pi y;\ \cos\pi x\,\cos\pi y\bigr).
\end{equation}
Solo aparece la diferencia $D_{11}-D_{12}=2\mu$. El coeficiente de corte $D_{33}$ y la parte
$\lambda$ común a $D_{11}$ y $D_{12}$ desaparecen: no influyen ni en $\bm b$ ni en $\uM$.

\textbf{Consecuencia.} Si el código tuviera un error en $D_{33}$ o en $\lambda$, el problema que
resolvería en realidad (con la misma $\bm b$ y los mismos valores de contorno) seguiría teniendo a
$\uM$ como solución exacta, porque las ecuaciones se cumplen con cualquier $D_{33}$ y cualquier
$\lambda$, y la solución es única. El MEF convergería a $\uM$ con los órdenes teóricos y la prueba
pasaría. Hay una segunda consecuencia: las matrices de tensión plana y de deformación plana tienen
el mismo $\mu$ y difieren solo en $\lambda$ ($\dato{lam}$ frente a $\dato{lamDP}$ con estos datos),
así que esta solución no distingue un estado plano del otro. En particular, el término fuente es el
mismo en ambos.

La tabla~\ref{tab:sabotajeA} lo comprueba sobre el propio EduFEM: se introdujeron a propósito
errores en \cod{fem/constitutive.py} y se corrieron las cuatro configuraciones del MMS de la tesis
con su criterio (\cod{experimento_sabotaje.py}).

\begin{table}[H]
\centering\small
\caption{Errores introducidos en la matriz $\bm D$ de EduFEM y resultado del MMS de la tesis. Órdenes
observados entre $N=16$ y $N=32$ (rango en las cuatro configuraciones).}\label{tab:sabotajeA}
\setlength{\tabcolsep}{5pt}
\begin{tabular}{@{}lccccc@{}}
\toprule
 & \multicolumn{2}{c}{Q4} & \multicolumn{2}{c}{Q9} & \\
\cmidrule(lr){2-3}\cmidrule(lr){4-5}
Error introducido & $L^2$ & $H^1$ & $L^2$ & $H^1$ & Pasan\\
\midrule
\tabla{sabotajeA}
\bottomrule
\end{tabular}
\end{table}

\begin{itemize}[leftmargin=1.2em]
  \item Si la rama de deformación plana devolviera la matriz de tensión plana, las ocho pruebas
    pasarían: las configuraciones de deformación plana darían exactamente las cifras de las de
    tensión plana. Ese error sí lo detecta otro test del repositorio,
    \cod{tests/test_vv_extensions.py} (prueba 8 de 8, von Mises en deformación plana), con un
    error del $\dato{vmsa}$\,\%.
  \item Si $D_{33}$ estuviera duplicado ---la confusión clásica entre la deformación angular
    $\gamma_{xy}$ y la componente tensorial $\varepsilon_{xy}=\gamma_{xy}/2$---, también pasarían
    las ocho. Lo detecta la membrana de Cook: con Q9 de $16\times16$, $u_y = \dato{cooksb}$ frente a
    $\dato{cookorig}$ sin el error.
  \item El control confirma que la prueba funciona para lo que sí ejercita: si cambia $D_{11}-D_{12}$,
    el código converge a otra solución, el error deja de bajar y el orden observado cae a cero.
\end{itemize}

Salari y Knupp describen el mismo fenómeno en uno de sus ejemplos: una malla curvilínea ortogonal
«no fue una buena elección, porque los términos cruzados de la métrica son nulos y, si tuvieran
errores, no aparecerían en la prueba de convergencia»~\cite[p.~69]{salari}.

\subsection*{Cómo elegir una solución que ejercite todos los términos}

La solución sirve si $D_{11}$, $D_{12}$ y $D_{33}$ entran en $\bm b$ de forma independiente. Una
opción es conservar $u_M$ y cambiar la frecuencia de $v_M$ en $y$:
\[
  u_M = \sin\pi x\,\sin\pi y, \qquad v_M = \cos\pi x\,\cos 2\pi y .
\]
Con la matriz genérica, el mismo cálculo de la sección~\ref{sec:paso3} da
\begin{equation}\label{eq:balternativa}
\begin{aligned}
  b_x &= \pi^2\bigl[(D_{11}+D_{33})\,\sin\pi x\,\sin\pi y - 2\,(D_{12}+D_{33})\,\sin\pi x\,\sin 2\pi y\bigr],\\
  b_y &= \pi^2\bigl[(4D_{11}+D_{33})\,\cos\pi x\,\cos 2\pi y - (D_{12}+D_{33})\,\cos\pi x\,\cos\pi y\bigr].
\end{aligned}
\end{equation}
Aparecen tres combinaciones: $D_{11}+D_{33}$, $D_{12}+D_{33}$ y $4D_{11}+D_{33}$. Son linealmente
independientes (el determinante de sus coeficientes vale $-3$), así que cualquier error en uno de
los tres coeficientes cambia $\bm b$. La frecuencia distinta en $y$ es necesaria: con la misma
frecuencia en $x$ y en $y$ se cumple $\partial^2 v/\partial x^2 = \partial^2 v/\partial y^2$, y
$D_{33}$ solo aparece sumado a $D_{11}$ o a $D_{12}$. La tabla~\ref{tab:sabotajeB} repite el
experimento con esta solución, en tensión plana (TP) y en deformación plana (DP).

\begin{table}[H]
\centering\small
\caption{Los mismos errores con la solución alternativa: orden observado $L^2$ / $H^1$ y veredicto.}\label{tab:sabotajeB}
\setlength{\tabcolsep}{4.5pt}
\begin{tabular}{@{}lcccc@{}}
\toprule
Error introducido & TP, Q4 & TP, Q9 & DP, Q4 & DP, Q9\\
\midrule
\tabla{sabotajeB}
\bottomrule
\end{tabular}
\end{table}

Con la solución alternativa, el código original pasa las cuatro pruebas. Con la matriz de
deformación plana cambiada por la de tensión plana fallan las dos de deformación plana (las de
tensión plana pasan, como corresponde, porque ahí no hay error), y con $D_{33}$ duplicado fallan
las cuatro.

\begin{aviso}{Qué significa para la tesis}
La matriz $\bm D$ de EduFEM es correcta: coincide con la fórmula~\eqref{eq:Dnum} y la
cubren otros tests. Lo que cambia es qué verifica el MMS tal como está programado: el ensamblaje,
la fuerza de cuerpo, las restricciones no homogéneas, el mapeo de elementos distorsionados y la
combinación $2\mu$ de la matriz $\bm D$, pero no $\lambda$ ni $D_{33}$. Por lo tanto, sus
configuraciones de deformación plana no verifican por sí solas la matriz de deformación plana.
\end{aviso}

% ==================================================================================
\section{Resumen del procedimiento}\label{sec:resumen}

Para repetir el método con otra solución, otro elemento u otro código:

\begin{enumerate}[leftmargin=1.6em]
  \item Escribir las ecuaciones que resuelve el código y su orden teórico ($p+1$ en $L^2$ y $p$ en
    $H^1$ para funciones de forma de grado completo $p$).
  \item Diseñar la prueba: dominio, material, elemento y tipo de restricción. Conviene empezar con
    restricciones en todo el contorno.
  \item Elegir $\uM$: suave, que no pertenezca al espacio de los elementos, con valores no nulos en
    el contorno y que ejercite todos los términos. Esto último se comprueba sustituyendo una matriz
    $\bm D$ genérica y viendo que cada coeficiente sobrevive en $\bm b$.
  \item Deducir $\bm b = -\nabla\cdot\bm\sigma(\uM)$ y los valores de contorno, entregarlos al código
    y obtener $\bm u_h$.
  \item Calcular $\|\eh\|_{L^2}$ y $|\eh|_{H^1}$ con una cuadratura de un orden más que la rigidez.
  \item Refinar la malla ($N$, $2N$, $4N$, \dots) y calcular $p_{\text{obs}} = \log_2(e_N/e_{2N})$.
  \item Comparar con el orden teórico en el par más fino: la prueba pasa si
    $|p_{\text{obs}} - p| < 0,5$.
  \item Si no pasa, revisar primero la prueba (la deducción de $\bm b$, los valores de contorno, el
    cálculo de las normas y si las mallas ya están en el rango asintótico) y solo después buscar el
    error en el código~\cite[pp.~33--35]{salari}.
\end{enumerate}

% ==================================================================================
\appendix
\section{Correspondencia con el código de EduFEM}\label{anx:codigo}

\begin{table}[H]
\centering\small
\caption{Dónde hace EduFEM cada paso del ejemplo.}\label{tab:codigo}
\renewcommand{\arraystretch}{1.15}
\begin{tabular}{@{}>{\raggedright\arraybackslash}p{0.34\textwidth}>{\raggedright\arraybackslash}p{0.60\textwidth}@{}}
\toprule
Paso & Archivo y función\\
\midrule
Solución manufacturada y su gradiente & \cod{tests/vv_mms.py}: \cod{u_M}, \cod{grad_u_M}\\
Fuerza de cuerpo~\eqref{eq:b} & \cod{tests/vv_mms.py}: \cod{make_body_force_fn}\\
Malla $N\times N$ & \cod{models/mesh_utils.py}: \cod{generate_structured_quad_mesh}\\
Nodos del contorno & \cod{tests/vv_mms.py}: \cod{boundary_nodes_topological}\\
Restricciones no homogéneas & \cod{models/project.py}: \cod{ProjectModel.set_boundary_condition}, con \texttt{ux} y \texttt{uy} no nulos\\
Rigidez~\eqref{eq:ke} y ensamblaje & \cod{fem/batch.py}: \cod{stiffness_batch}, \cod{assemble_sparse}; versión legible en \cod{fem/stiffness.py}\\
Fuerzas de la fuerza de cuerpo & \cod{fem/assembly.py}: \cod{_body_force_at_gauss_points}; \cod{fem/batch.py}: \cod{body_force_batch}\\
Sustitución estática~\eqref{eq:reducido} & \cod{fem/solver.py}: \cod{apply_boundary_conditions}\\
Normas~\eqref{eq:normas} & \cod{fem/error_norms.py}: \cod{compute_error_norms}\\
Tensiones recuperadas $\bm\sigma^*$ & \cod{fem/error_norms.py}: \cod{compute_stress_recovery_error}\\
Orden observado y criterio & \cod{tests/vv_mms.py}: \cod{compute_rates}, \cod{main}\\
\bottomrule
\end{tabular}
\end{table}

\section{La réplica independiente}\label{anx:replica}

El guion \cod{mms_independiente.py}, en la misma carpeta que este documento, repite el paso 4 y el
refinamiento sin importar nada de EduFEM. Usa solo NumPy y SciPy, recorre los elementos uno por
uno y está escrito para leerse. Se ejecuta con \texttt{python mms\_independiente.py} (unos 20~s) e
imprime el desplazamiento del nodo central para $N=2$ y la tabla~\ref{tab:refinamiento} con sus
órdenes. Internamente numera los nodos de cada elemento por filas (producto tensorial de dos
bases de Lagrange 1D), no en sentido antihorario como la sección~\ref{sec:mano}; el resultado no
depende de esa numeración.

Los otros dos guiones de la carpeta regeneran el documento. \cod{generar_datos.py} calcula todas
las cifras (escribe \cod{datos.tex}, \cod{convergencia.dat} y \cod{pendientes.dat}) y las contrasta
con EduFEM. \cod{experimento_sabotaje.py} corre los experimentos de la sección~\ref{sec:limites}
(escribe \cod{datos_sabotaje.tex}, unos 3~minutos).

\lstinputlisting{mms_independiente.py}

% ==================================================================================
\begin{thebibliography}{9}\small
\bibitem{salari} Salari K, Knupp P. Code verification by the method of manufactured solutions.
  Albuquerque (NM): Sandia National Laboratories; 2000. Informe n.º SAND2000-1444.
\bibitem{hughes} Hughes TJR. The finite element method: linear static and dynamic finite element
  analysis. Englewood Cliffs (NJ): Prentice-Hall; 1987.
\bibitem{strang} Strang G, Fix GJ. An analysis of the finite element method. Englewood Cliffs (NJ):
  Prentice-Hall; 1973.
\end{thebibliography}

\end{document}
